\subsection{1644—1652：越过关隘，南方仍在移动}

顺治元年四月，山海关把大顺、关宁军和清摄政政权的危机压在同一处关隘。李自成进入北京后，关宁军与大顺在兵饷、统属和互信上已经破裂；多尔衮利用吴三桂的处境，使八旗获得越关的行动空间。会战打开的是华北道路，并不是天下已经归属。五月入北京后，清军接收的首先是仓储、官署、粮道和观望中的官员；九月顺治帝自盛京抵京，皇帝驻跸、部院和礼制中心才开始汇合。冬季的安民、招纳、征发和整顿命令，显示一支行军军队正在试图成为能收发文书的朝廷。

这条线很快在江南折断成多条战线。1645年南京失守、弘光帝被俘，清廷又重申薙发易服；命令把发式当作政治归属的可见标记，却在扬州、江阴等地同守城、逃亡、征发和杀戮纠缠。城陷、围城和严重暴力是征服史不能略去的事实；《扬州十日记》《江阴城守纪》、地方笔记与方志所保存的日数、伤亡和见闻，则不能被改写成全境的精确统计。弘光的覆亡也没有使南方同时结束：鲁王集团、隆武、绍武和永历的政治中心仍在浙闽、两广和西南移动。

1646年至1650年，清军的攻取、南明朝廷的流转和将领的反复选择交织在一起。隆武、绍武政权先后覆亡，永历名义在两广和西南继续被不同军队借用；1648年金声桓、李成栋改奉南明，次年南昌再失、1650年广州再陷，又说明先前的“接收”并未冻结地方的选择。城池、道路、饷源和港口并不以同一种速度改易归属。清方军报能够钉住其所报的攻守与处分；南明奏疏、地方记事和城守记录保留的是败亡朝廷与城市共同体的另一种时间，二者都不能代替整个战区。

顺治七年多尔衮去世，次年顺治帝亲政，摄政枢纽随之改组；皇帝名义的加强却仍要经过宗室、旗务、满汉官僚和前线军饷才能抵达地方。1652年李定国在桂林、衡州一带的战事迫使清廷重新估量湘桂道路与机动兵力，西南的胜负并未随北京的亲政而预先决定。到这一年，读者所能看见的不是一条完成的疆界，而是仪式、降附、逃散和再选择不断彼此牵动。

\paragraph{材料位置。} 《清世祖实录》、内国史院满文档案和题本记录新政权的命令、接收与军务；南明奏疏、《永历实录》及地方城守材料提供对手和地方的时间线。它们可以相互校正事件顺序，却不能让求援文书原貌、各军兵额、攻城伤亡或每一处执行情形获得并不存在的精确性。

\section{从关门到京师}

顺治元年四月，山海关不是一扇一经开启便可直达天下的门，而是一处把几种危机压在一起的关隘。李自成入北京之后，关宁军同大顺之间的兵饷、统属与互信已经破裂；多尔衮则把这一裂缝变成八旗越过关隘的机会。清军与吴三桂部在山海关击败大顺军，因而取得进入华北的通道。这一结果可以立为编年事实；但若把后来流行的“乞师”故事逐句当作战前文件，把各军人数、会师次序和战场上的每一次转折都写成定论，证据便承受不住这样的重压。

《清世祖实录》中的军报把胜利纳入摄政政权的行动叙述，清初内国史院保存的满文关宁军务档案则让人看见命令、呈报和接收如何在文书链中移动。二者首先是清廷留下的行政记录：它们能证明朝廷收到什么消息、以何种语句处分军政，却不能自动代替关内各军的全部经验。后出的《清史稿》适合用来互校官名和大事顺序，但它是民国初年未定稿的官修史书，不是顺治四月的现场记录。吴三桂求援文书的原貌、各方兵额以及战场过程，仍须保留为由清廷叙事和后出文本共同塑形的问题。

五月入北京，也不应被压缩成“占领京城”四字。大顺撤离、明廷解体之后，城中留下的是仓储、官署、粮饷、逃散的人群和等待选择的官员。多尔衮率军入城，接收的是明朝京师的行政与礼制空间；到冬季，摄政王又以安民、征发、招纳明官和整顿官署的命令，尝试把军队占有的城池变成可以传递文书的治所。诏令的发布是可靠的，题本中可见的交接也证明部分官署重新运转；可是受任者是否到任、一纸命令是否越过城门和驿路、地方是否照办，不能从诏书的语气倒推出去。

北京地方记事和《甲申传信录》保存了另一种尺度：它们所见的是城内消息、秩序与恐惧，而非摄政政权所要展示的承接。它们同样有作者位置、传抄与见闻范围的限制，却正可防止以中央档案的沉默替代城市生活。顺治帝九月自盛京抵达北京、举行入城和定都仪式，标志着皇帝驻跸与部院中心向北京集中；盛京却仍保有满洲政治和象征上的重量。仪式建立的是持续的中央名义，并不等于各地已经承认，尤其不等于南方战场已经结束。

\section{摄政并非一条畅通的命令线}

多尔衮的摄政，是清初征服得以调动兵马、官员和名义资源的枢纽，但不是一架无阻运转的机器。京师需要粮饷，部院需要熟悉旧制的书吏和降官，前线又不断索要兵丁、漕运和军需。清廷档案最擅长呈现这条由奏报、批示、部议到任命的路径；地方档案、方志和南方对手的材料则不断提醒人们，路径在何处中断、被拖延，或被战事改写。把“招降”写成一个完成状态，会遮蔽降官仍须在家族、旧部、军饷和地方安全之间反复选择的处境。

顺治七年十二月，多尔衮在古北口外行猎途中去世。其死期与丧仪可以由《清世祖实录》和内国史院的遗命、丧仪档相互钉住；死因以及死后政治清算的动机，却不能只照后来定性的语言复述。次年顺治帝亲政，摄政安排撤除，皇帝周边的任官与决策网络重新组合。然而，署在皇帝名下的命令仍要依靠宗室、旗务、满汉官僚与南方将领落实。摄政结束，不是中央能力忽然圆满的时刻。

这一点在顺治帝去世后更为清楚。顺治十八年，八岁的玄烨即位，索尼、苏克萨哈、遏必隆、鳌拜辅政。继位与四辅政大臣的名单见于《清世祖实录》、起居注和遗诏档；至康熙四年鳌拜影响上升、康熙六年索尼去世、康熙八年鳌拜被拘捕，则应分别阅读《清圣祖实录》、起居注与内阁案卷。尤其不能把1669年的结果倒灌回此前：后来的鳌拜案和人物传记最容易把本来由议政、旗务、人事和军费共同构成的辅政政治，写成一人独揽权力的必然前史。康熙拘捕鳌拜使皇帝能够更直接地统摄人事军政，却也仍须在既有制度与旗人贵族网络中行使这种权力。

\section{南京失守以后：城池、朝廷与身体}

清军南下的速度使京师接收和江南战争看似连成一线，实际却是一系列由水陆节点、粮道和地方选择构成的战区。顺治二年春夏之交，清军沿运河、淮扬推进，攻取南京；弘光帝朱由崧不久在芜湖一带被俘。清方的《清世祖实录》与《清史稿·世祖本纪》可以确立城陷、押解这一主干，南明弘光朝奏疏、《南明史料》所收文书和江宁地方记事则保留了另一面：朝廷如何失去护卫与指挥，官员和军队怎样逃散、观望或改易归属。前者是胜利政权收到的战报，后者多是败亡朝廷及城市中的急迫文字；两类材料不能互相注销，也不应被合成为一种全知叙述。

南京失去首都地位，并没有令“明亡”在南方成为一个同时发生的事件。弘光名义瓦解后，鲁王集团在浙东和海上活动，隆武朝廷在福建建立，绍武与永历又在两广相继争夺名义中心。清军取得杭州、绍兴等城，或攻入广州，并不自动说明一省的道路、岛屿、山地和县城都已归于同一控制。清廷军报常以府城、俘获和招降排列战果；《福建通志》《广东通志》及各地府县志可以补足部分地名和后续秩序，却大多成书较晚、也会沿袭胜败叙事。南明文集和《永历实录》把流动朝廷、将领联络与退守路线连成另一条线索，但它们也服务于保存正统、追责或追忆的需要。史实的难处不在于挑选哪一方“说真话”，而在于说明每种文本首先能够证明什么。

顺治二年六月重申的薙发易服命令，使政治归属进入身体的可见处。实录中的命令证明清廷以发式服制辨认新旧政治身份；兵科题本和江南府县志的相关条目，则可见命令在若干地方与征发、守城和治安纠缠在一起。它不能证明全国在同一日、以同一种方式服从。有人遵从，有人逃亡，有人把发式当作拒绝交付城池或粮饷的标识；不同府县的速度、强度与后果必须逐地追问。

因此，扬州、江阴和1650年的广州特别不宜用一个整齐的死亡数字盖棺。扬州失陷、江阴围城失守、广州再陷，以及攻城后发生严重杀戮、掳掠和损毁，都是必须写入征服史的事实。可是清军战报所陈的战果，不是伤亡统计；《扬州十日记》《江阴城守纪》、广州地方笔记和后出的府县志，保存了幸存者、编纂者或地方共同体的恐惧与纪念，也各有成书、传抄、修辞和数字来源的问题。“十日”、围城天数和死亡人数不应被伪装成精确普查，更不能由一部文本推算江南或广东整体的死亡。较为诚实的写法是同时承认暴力及其长期记忆，说明现有材料不能可靠合算总数，并把不同数字标回提出它的作者、地点与文类。

\section{南方并未在一役中归入清廷}

1646年隆武帝在闽北被俘杀、绍武朝在广州结束之后，朱由榔在肇庆即位，永历年号成为两广和西南多支力量可以借用的名义。清廷军报将这看作继续清剿的对象；《永历实录》、永历朝奏疏和地方志所呈现的，却是一个不断移动、依赖地方军队、饷源与山海通道的朝廷。1648年金声桓在南昌、李成栋在广东改奉南明，说明此前的“接收”未能冻结将领与地方的选择。1649年南昌被攻取、1650年广州再陷固然削弱了这些反转的中心，却不能把战后的税源、流民和港口恢复一并写成已经解决。

永历六年，李定国攻取桂林、在衡州击败清军，迫使清廷重新估量湘桂道路、军需和机动兵力。双方的兵数、死伤和将领死因在清廷军报与南明记录之间多有不同；可以肯定战场结果，不能替任何一方把报捷数字升级为无争议的计数。此后孙可望与永历朝的冲突及其降清，削弱了南明的军事协作；清军自贵州、四川等方向入滇，永历朝廷离开昆明，终至进入缅甸。这里又有一层不能省略的证据边界：清方材料擅长记录追捕与军政调度，缅甸宫廷的《琉璃宫史》传统和欧洲旅行记所保留的接待、路线与交涉，并不必然与之重合。把缅甸只写作清军追击地图上的空白，既抹掉了东吁王朝的选择，也误解了流亡朝廷的补给条件。

1659年郑成功北上围攻南京而后撤退，显示长江下游并非已无战事；翌年清军经营云南城镇和驿路，也说明占有昆明不等于控制边地。1662年，朱由榔由缅甸被押回后处死，南明皇帝的名义中心至此终结。中缅材料对于移交过程、处死地点和责任分配仍须并读。夔东抗清据点直到1664年才被清军攻破，郑氏又在海上和台湾维持政权；这些后续并非枝节，而是对“1644年入关即全国统一”这一方便说法的校正。

从山海关到北京，从南京到西南边地，清初形成的不是一条单向延伸的疆界，而是军队、文书、仪式、降附、逃亡与再选择彼此牵动的过程。清廷档案使我们能够看见命令如何被制作和宣告；官修史书为后来者编排出帝国的时间；地方材料留下城池和社区承受的断片；南明叙事则保存失败者仍在行动的政治地理。把它们放在同一页上，不是为了把分歧磨平，而是为了让征服、摄政与整合都回到它们当时尚未完成的样子。


\subsection{事件链与材料限度}
关隘、京城与官署的材料并不构成一条无摩擦的接收线。以下事件保留各自的行动者、空间、后果与证据限度。

\paragraph{顺治元年五月（1644年6月前后）：多尔衮率军进入北京并接收明朝京城的行政与礼制空间}
\eventanchor{EQ-1644-BEIJING-ENTRY}{多尔衮率军进入北京并接收明朝京城的行政与礼制空间}。本节点叙述该事件本身。清摄政政权与北京在京师与多地军政线路相接的尺度的处境首先受大顺撤离与明廷崩溃留下军政真空，清军需获取京师粮饷、官僚和象征中心制约。《清世祖实录》顺治元年五月；《清史稿·世祖本纪》；《甲申传信录》及北京地方记事留下可核的线索；可辨的后果是清廷以北京为治所，开始将军事占领转化为对明朝制度资源的继承和重组。原有置信注记：入城与接收可靠；城内秩序、迎降者范围和掠夺状况须以不同见证互校

\paragraph{顺治元年五月（1644年6月前后）：多尔衮率军进入北京并接收明朝京城的行政与礼制空间}
\eventanchor{EQ-1644-BEIJING-ENTRY-AFTERMATH}{多尔衮率军进入北京并接收明朝京城的行政与礼制空间} 置于京师与多地军政线路相接的尺度中阅读，本节点追踪「多尔衮率军进入北京并接收明朝京城的行政与礼制空间」之后可见的余波。大顺撤离与明廷崩溃留下军政真空，清军需获取京师粮饷、官僚和象征中心构成行动者清摄政政权与北京必须面对的条件。取证从《清世祖实录》顺治元年五月；《清史稿·世祖本纪》；《甲申传信录》及北京地方记事出发，因而只能把变化谨慎地连到清廷以北京为治所，开始将军事占领转化为对明朝制度资源的继承和重组。原有置信注记：入城与接收可靠；城内秩序、迎降者范围和掠夺状况须以不同见证互校

\paragraph{顺治元年五月（1644年6月前后）：多尔衮率军进入北京并接收明朝京城的行政与礼制空间}
\eventanchor{EQ-1644-BEIJING-ENTRY-DECISION}{多尔衮率军进入北京并接收明朝京城的行政与礼制空间} 标记本节点定位围绕「多尔衮率军进入北京并接收明朝京城的行政与礼制空间」形成的处置与调度记录，涉及清摄政政权与北京。京师与多地军政线路相接的尺度不是抽象背景：大顺撤离与明廷崩溃留下军政真空，清军需获取京师粮饷、官僚和象征中心。《清世祖实录》顺治元年五月；《清史稿·世祖本纪》；《甲申传信录》及北京地方记事所见的顺序随后指向清廷以北京为治所，开始将军事占领转化为对明朝制度资源的继承和重组，但原有置信注记：入城与接收可靠；城内秩序、迎降者范围和掠夺状况须以不同见证互校

\paragraph{顺治元年五月（1644年6月前后）：多尔衮率军进入北京并接收明朝京城的行政与礼制空间}
在京师与多地军政线路相接的尺度，\eventanchor{EQ-1644-BEIJING-ENTRY-PRELUDE}{多尔衮率军进入北京并接收明朝京城的行政与礼制空间} 让清摄政政权与北京的一个选择或压力有了可查的坐标。本节点回看「多尔衮率军进入北京并接收明朝京城的行政与礼制空间」之前已可见的条件。前因可写为大顺撤离与明廷崩溃留下军政真空，清军需获取京师粮饷、官僚和象征中心；《清世祖实录》顺治元年五月；《清史稿·世祖本纪》；《甲申传信录》及北京地方记事限定了记录，后续变化是清廷以北京为治所，开始将军事占领转化为对明朝制度资源的继承和重组。原有置信注记：入城与接收可靠；城内秩序、迎降者范围和掠夺状况须以不同见证互校

\paragraph{顺治元年冬（1644年）：多尔衮以摄政王名义整顿京师官署、招降明官并发布安民与征发命令}
\eventanchor{EQ-1644-REGENT-ORDERS}{多尔衮以摄政王名义整顿京师官署、招降明官并发布安民与征发命令} 所关联的是清摄政政权在京师与多地军政线路相接的尺度中的行动。本节点叙述该事件本身。它从新政权需要恢复行政传递、维持京师粮饷并吸纳明代文官技术展开，借《清世祖实录》顺治元年十月至十二月；《清史稿·职官志》；中国第一历史档案馆藏内阁大库顺治元年题本〔京师官署交接〕进入可检验的年序，并以中央可用的部院和任官网络开始恢复，但忠诚、财力与军纪继续限制覆盖面影响下一段安排。原有置信注记：命令发布可靠；接受任官和地方执行不能从诏令直接推出

\paragraph{顺治元年冬（1644年）：多尔衮以摄政王名义整顿京师官署、招降明官并发布安民与征发命令}
京师与多地军政线路相接的尺度中的\eventanchor{EQ-1644-REGENT-ORDERS-AFTERMATH}{多尔衮以摄政王名义整顿京师官署、招降明官并发布安民与征发命令} 不能离开清摄政政权来解释。本节点追踪「多尔衮以摄政王名义整顿京师官署、招降明官并发布安民与征发命令」之后可见的余波。新政权需要恢复行政传递、维持京师粮饷并吸纳明代文官技术说明其形成的条件；《清世祖实录》顺治元年十月至十二月；《清史稿·职官志》；中国第一历史档案馆藏内阁大库顺治元年题本〔京师官署交接〕提供来源路径，中央可用的部院和任官网络开始恢复，但忠诚、财力与军纪继续限制覆盖面则是材料能够支持的近时结果。原有置信注记：命令发布可靠；接受任官和地方执行不能从诏令直接推出

\paragraph{顺治元年冬（1644年）：多尔衮以摄政王名义整顿京师官署、招降明官并发布安民与征发命令}
\eventanchor{EQ-1644-REGENT-ORDERS-DECISION}{多尔衮以摄政王名义整顿京师官署、招降明官并发布安民与征发命令} 把清摄政政权放回京师与多地军政线路相接的尺度的道路、城镇、边界或制度网络。本节点定位围绕「多尔衮以摄政王名义整顿京师官署、招降明官并发布安民与征发命令」形成的处置与调度记录。新政权需要恢复行政传递、维持京师粮饷并吸纳明代文官技术。本段采用《清世祖实录》顺治元年十月至十二月；《清史稿·职官志》；中国第一历史档案馆藏内阁大库顺治元年题本〔京师官署交接〕核对其顺序，随后可追到中央可用的部院和任官网络开始恢复，但忠诚、财力与军纪继续限制覆盖面。原有置信注记：命令发布可靠；接受任官和地方执行不能从诏令直接推出

\paragraph{顺治元年冬（1644年）：多尔衮以摄政王名义整顿京师官署、招降明官并发布安民与征发命令}
本节点回看「多尔衮以摄政王名义整顿京师官署、招降明官并发布安民与征发命令」之前已可见的条件，故\eventanchor{EQ-1644-REGENT-ORDERS-PRELUDE}{多尔衮以摄政王名义整顿京师官署、招降明官并发布安民与征发命令} 不只是一个年份标签。清摄政政权在京师与多地军政线路相接的尺度承受的压力来自新政权需要恢复行政传递、维持京师粮饷并吸纳明代文官技术。《清世祖实录》顺治元年十月至十二月；《清史稿·职官志》；中国第一历史档案馆藏内阁大库顺治元年题本〔京师官署交接〕可回查该节点；中央可用的部院和任官网络开始恢复，但忠诚、财力与军纪继续限制覆盖面是其进入后续年序的方式。原有置信注记：命令发布可靠；接受任官和地方执行不能从诏令直接推出

\paragraph{顺治元年四月（1644年5月，换算待核）：清军与吴三桂在山海关击败李自成军，打开进入华北的军事通道}
\eventanchor{EQ-1644-SHANHAIGUAN-AFTERMATH}{清军与吴三桂在山海关击败李自成军，打开进入华北的军事通道} 所见的行动者为清摄政政权、吴三桂与大顺，空间尺度为京师与多地军政线路相接的尺度。本节点追踪「清军与吴三桂在山海关击败李自成军，打开进入华北的军事通道」之后可见的余波。李自成入北京后关宁军与大顺在兵饷、统属和互信上破裂，多尔衮把危机转为入关机会使它成为具体事件链的一环；《清世祖实录》顺治元年四月山海关军报；《清史稿·世祖本纪》；《清初内国史院满文档案》顺治元年四月关宁军务档与清军获得穿越关隘的行动空间，北京政权争夺转为华北与南方的多战区征服。原有置信注记：核心战事可靠；吴三桂求援文书原貌、各军兵数和战场过程受清廷与后出叙事塑形

\paragraph{顺治元年四月（1644年5月，换算待核）：清军与吴三桂在山海关击败李自成军，打开进入华北的军事通道}
从京师与多地军政线路相接的尺度看，\eventanchor{EQ-1644-SHANHAIGUAN-DECISION}{清军与吴三桂在山海关击败李自成军，打开进入华北的军事通道} 留下本节点定位围绕「清军与吴三桂在山海关击败李自成军，打开进入华北的军事通道」形成的处置与调度记录的材料坐标。清摄政政权、吴三桂与大顺所面对的条件是李自成入北京后关宁军与大顺在兵饷、统属和互信上破裂，多尔衮把危机转为入关机会。《清世祖实录》顺治元年四月山海关军报；《清史稿·世祖本纪》；《清初内国史院满文档案》顺治元年四月关宁军务档支持的结果为清军获得穿越关隘的行动空间，北京政权争夺转为华北与南方的多战区征服。原有置信注记：核心战事可靠；吴三桂求援文书原貌、各军兵数和战场过程受清廷与后出叙事塑形

\paragraph{顺治元年四月（1644年5月，换算待核）：清军与吴三桂在山海关击败李自成军，打开进入华北的军事通道}
\eventanchor{EQ-1644-SHANHAIGUAN-PRELUDE}{清军与吴三桂在山海关击败李自成军，打开进入华北的军事通道} 记录清摄政政权、吴三桂与大顺在京师与多地军政线路相接的尺度的一次变化。本节点回看「清军与吴三桂在山海关击败李自成军，打开进入华北的军事通道」之前已可见的条件。其发生与李自成入北京后关宁军与大顺在兵饷、统属和互信上破裂，多尔衮把危机转为入关机会相连；《清世祖实录》顺治元年四月山海关军报；《清史稿·世祖本纪》；《清初内国史院满文档案》顺治元年四月关宁军务档把事件系回来源，清军获得穿越关隘的行动空间，北京政权争夺转为华北与南方的多战区征服显示压力如何向后转移。原有置信注记：核心战事可靠；吴三桂求援文书原貌、各军兵数和战场过程受清廷与后出叙事塑形

\paragraph{顺治元年九月（1644年10月前后）：顺治帝自盛京抵北京，清廷举行定都与皇帝入城仪式}
\eventanchor{EQ-1644-SHUNZHI-BEIJING}{顺治帝自盛京抵北京，清廷举行定都与皇帝入城仪式} 的地点与行动范围属于京师与多地军政线路相接的尺度，行动者是清。本节点叙述该事件本身。摄政政权虽占北京，仍需皇帝亲临与礼制重置来建立持续中央名义。读《清世祖实录》顺治元年九月；《清史稿·世祖本纪》；《清初内国史院满文档案》北京迁都礼仪档时可见其记录边界，最小后果只能写作北京成为皇帝驻跸与部院中心，盛京仍保有满洲政治和象征地位。原有置信注记：日期与仪式主干可靠；官修礼仪记录不能等同各地承认

\paragraph{顺治元年九月（1644年10月前后）：顺治帝自盛京抵北京，清廷举行定都与皇帝入城仪式}
\eventanchor{EQ-1644-SHUNZHI-BEIJING-AFTERMATH}{顺治帝自盛京抵北京，清廷举行定都与皇帝入城仪式} 是清在京师与多地军政线路相接的尺度留下的编年标记。本节点追踪「顺治帝自盛京抵北京，清廷举行定都与皇帝入城仪式」之后可见的余波。它从摄政政权虽占北京，仍需皇帝亲临与礼制重置来建立持续中央名义走来；《清世祖实录》顺治元年九月；《清史稿·世祖本纪》；《清初内国史院满文档案》北京迁都礼仪档固定可证部分，北京成为皇帝驻跸与部院中心，盛京仍保有满洲政治和象征地位构成下一步的可见结果。原有置信注记：日期与仪式主干可靠；官修礼仪记录不能等同各地承认

\paragraph{顺治元年九月（1644年10月前后）：顺治帝自盛京抵北京，清廷举行定都与皇帝入城仪式}
本节点定位围绕「顺治帝自盛京抵北京，清廷举行定都与皇帝入城仪式」形成的处置与调度记录：\eventanchor{EQ-1644-SHUNZHI-BEIJING-DECISION}{顺治帝自盛京抵北京，清廷举行定都与皇帝入城仪式}。清在京师与多地军政线路相接的尺度的处境可由摄政政权虽占北京，仍需皇帝亲临与礼制重置来建立持续中央名义说明。《清世祖实录》顺治元年九月；《清史稿·世祖本纪》；《清初内国史院满文档案》北京迁都礼仪档是本段材料入口，因而只能据此追踪到北京成为皇帝驻跸与部院中心，盛京仍保有满洲政治和象征地位。原有置信注记：日期与仪式主干可靠；官修礼仪记录不能等同各地承认

\paragraph{顺治元年九月（1644年10月前后）：顺治帝自盛京抵北京，清廷举行定都与皇帝入城仪式}
\eventanchor{EQ-1644-SHUNZHI-BEIJING-PRELUDE}{顺治帝自盛京抵北京，清廷举行定都与皇帝入城仪式} 发生在京师与多地军政线路相接的尺度，牵动清。本节点回看「顺治帝自盛京抵北京，清廷举行定都与皇帝入城仪式」之前已可见的条件。摄政政权虽占北京，仍需皇帝亲临与礼制重置来建立持续中央名义是其结构性条件；《清世祖实录》顺治元年九月；《清史稿·世祖本纪》；《清初内国史院满文档案》北京迁都礼仪档留下的证据把读者带到北京成为皇帝驻跸与部院中心，盛京仍保有满洲政治和象征地位。原有置信注记：日期与仪式主干可靠；官修礼仪记录不能等同各地承认



\subsection{事件链与材料限度}
南方战场的城池、江路与地方社会没有同速改变。下列节点保存征服、迁徙与接收之间的材料距离。

\paragraph{顺治二年六月（1645年）：清廷重申薙发易服命令，在江南等地引发服从、逃亡与抵抗}
\eventanchor{EQ-1645-HAIR-ORDER}{清廷重申薙发易服命令，在江南等地引发服从、逃亡与抵抗}。本节点叙述该事件本身。清与江南各地在省、府县与地方社会相互牵动的尺度的处境首先受清廷以发式服制识别新旧政治归属，军事推进又使命令与治安相互捆绑制约。《清世祖实录》顺治二年六月；《清史稿·舆服志》；内阁大库顺治二年兵科题本与江南府县志薙发条留下可核的线索；可辨的后果是服制成为忠降识别和集体动员的触发点，若干江南城镇的冲突加剧征服成本。原有置信注记：命令文本可靠；各地执行速度和社会反应高度不一，不能以诏书概括全国

\paragraph{顺治二年六月（1645年）：清廷重申薙发易服命令，在江南等地引发服从、逃亡与抵抗}
\eventanchor{EQ-1645-HAIR-ORDER-AFTERMATH}{清廷重申薙发易服命令，在江南等地引发服从、逃亡与抵抗} 置于省、府县与地方社会相互牵动的尺度中阅读，本节点追踪「清廷重申薙发易服命令，在江南等地引发服从、逃亡与抵抗」之后可见的余波。清廷以发式服制识别新旧政治归属，军事推进又使命令与治安相互捆绑构成行动者清与江南各地必须面对的条件。取证从《清世祖实录》顺治二年六月；《清史稿·舆服志》；内阁大库顺治二年兵科题本与江南府县志薙发条出发，因而只能把变化谨慎地连到服制成为忠降识别和集体动员的触发点，若干江南城镇的冲突加剧征服成本。原有置信注记：命令文本可靠；各地执行速度和社会反应高度不一，不能以诏书概括全国

\paragraph{顺治二年六月（1645年）：清廷重申薙发易服命令，在江南等地引发服从、逃亡与抵抗}
\eventanchor{EQ-1645-HAIR-ORDER-DECISION}{清廷重申薙发易服命令，在江南等地引发服从、逃亡与抵抗} 标记本节点定位围绕「清廷重申薙发易服命令，在江南等地引发服从、逃亡与抵抗」形成的处置与调度记录，涉及清与江南各地。省、府县与地方社会相互牵动的尺度不是抽象背景：清廷以发式服制识别新旧政治归属，军事推进又使命令与治安相互捆绑。《清世祖实录》顺治二年六月；《清史稿·舆服志》；内阁大库顺治二年兵科题本与江南府县志薙发条所见的顺序随后指向服制成为忠降识别和集体动员的触发点，若干江南城镇的冲突加剧征服成本，但原有置信注记：命令文本可靠；各地执行速度和社会反应高度不一，不能以诏书概括全国

\paragraph{顺治二年六月（1645年）：清廷重申薙发易服命令，在江南等地引发服从、逃亡与抵抗}
在省、府县与地方社会相互牵动的尺度，\eventanchor{EQ-1645-HAIR-ORDER-PRELUDE}{清廷重申薙发易服命令，在江南等地引发服从、逃亡与抵抗} 让清与江南各地的一个选择或压力有了可查的坐标。本节点回看「清廷重申薙发易服命令，在江南等地引发服从、逃亡与抵抗」之前已可见的条件。前因可写为清廷以发式服制识别新旧政治归属，军事推进又使命令与治安相互捆绑；《清世祖实录》顺治二年六月；《清史稿·舆服志》；内阁大库顺治二年兵科题本与江南府县志薙发条限定了记录，后续变化是服制成为忠降识别和集体动员的触发点，若干江南城镇的冲突加剧征服成本。原有置信注记：命令文本可靠；各地执行速度和社会反应高度不一，不能以诏书概括全国

\paragraph{顺治二年五月（1645年）：弘光帝朱由崧在芜湖一带被俘并被押往北京}
\eventanchor{EQ-1645-HONGGUANG-CAPTURE}{弘光帝朱由崧在芜湖一带被俘并被押往北京} 所关联的是清与南明弘光政权在京师与多地军政线路相接的尺度中的行动。本节点叙述该事件本身。它从南京失守后弘光朝没有持续的军事护卫与统一指挥，逃亡路线被清军和地方武装切断展开，借《清世祖实录》顺治二年五月；《清史稿·弘光帝本纪》；《明季南略》及江南地方志相关条目进入可检验的年序，并以以弘光名义维系的朝廷失能，隆武、永历政权在东南和西南重建名义中心影响下一段安排。原有置信注记：被俘及押解主干可靠；参与者责任和移交地点在纪传、地方记事间有异

\paragraph{顺治二年五月（1645年）：弘光帝朱由崧在芜湖一带被俘并被押往北京}
京师与多地军政线路相接的尺度中的\eventanchor{EQ-1645-HONGGUANG-CAPTURE-AFTERMATH}{弘光帝朱由崧在芜湖一带被俘并被押往北京} 不能离开清与南明弘光政权来解释。本节点追踪「弘光帝朱由崧在芜湖一带被俘并被押往北京」之后可见的余波。南京失守后弘光朝没有持续的军事护卫与统一指挥，逃亡路线被清军和地方武装切断说明其形成的条件；《清世祖实录》顺治二年五月；《清史稿·弘光帝本纪》；《明季南略》及江南地方志相关条目提供来源路径，以弘光名义维系的朝廷失能，隆武、永历政权在东南和西南重建名义中心则是材料能够支持的近时结果。原有置信注记：被俘及押解主干可靠；参与者责任和移交地点在纪传、地方记事间有异

\paragraph{顺治二年五月（1645年）：弘光帝朱由崧在芜湖一带被俘并被押往北京}
\eventanchor{EQ-1645-HONGGUANG-CAPTURE-DECISION}{弘光帝朱由崧在芜湖一带被俘并被押往北京} 把清与南明弘光政权放回京师与多地军政线路相接的尺度的道路、城镇、边界或制度网络。本节点定位围绕「弘光帝朱由崧在芜湖一带被俘并被押往北京」形成的处置与调度记录。南京失守后弘光朝没有持续的军事护卫与统一指挥，逃亡路线被清军和地方武装切断。本段采用《清世祖实录》顺治二年五月；《清史稿·弘光帝本纪》；《明季南略》及江南地方志相关条目核对其顺序，随后可追到以弘光名义维系的朝廷失能，隆武、永历政权在东南和西南重建名义中心。原有置信注记：被俘及押解主干可靠；参与者责任和移交地点在纪传、地方记事间有异

\paragraph{顺治二年五月（1645年）：弘光帝朱由崧在芜湖一带被俘并被押往北京}
本节点回看「弘光帝朱由崧在芜湖一带被俘并被押往北京」之前已可见的条件，故\eventanchor{EQ-1645-HONGGUANG-CAPTURE-PRELUDE}{弘光帝朱由崧在芜湖一带被俘并被押往北京} 不只是一个年份标签。清与南明弘光政权在京师与多地军政线路相接的尺度承受的压力来自南京失守后弘光朝没有持续的军事护卫与统一指挥，逃亡路线被清军和地方武装切断。《清世祖实录》顺治二年五月；《清史稿·弘光帝本纪》；《明季南略》及江南地方志相关条目可回查该节点；以弘光名义维系的朝廷失能，隆武、永历政权在东南和西南重建名义中心是其进入后续年序的方式。原有置信注记：被俘及押解主干可靠；参与者责任和移交地点在纪传、地方记事间有异

\paragraph{顺治二年闰六月至八月（1645年）：江阴在薙发命令和清军围攻中失守，地方记忆长期以守城与屠城叙述该事件}
\eventanchor{EQ-1645-JIANGYIN}{江阴在薙发命令和清军围攻中失守，地方记忆长期以守城与屠城叙述该事件} 所见的行动者为清军与江阴地方守军，空间尺度为省、府县与地方社会相互牵动的尺度。本节点叙述该事件本身。服制命令与守备危机把行政服从转成军事对抗，清军需确保沿江交通和粮道使它成为具体事件链的一环；《清世祖实录》顺治二年江南军报；《清史稿·世祖本纪》；《江阴县志》及《江阴城守纪》与清廷以军事压制恢复沿江节点；江阴成为地方忠义记忆，伤亡规模不能由单一叙事定量。原有置信注记：城陷主干可靠；围城天数、伤亡和人物言行多依赖地方后出文本

\paragraph{顺治二年闰六月至八月（1645年）：江阴在薙发命令和清军围攻中失守，地方记忆长期以守城与屠城叙述该事件}
从省、府县与地方社会相互牵动的尺度看，\eventanchor{EQ-1645-JIANGYIN-AFTERMATH}{江阴在薙发命令和清军围攻中失守，地方记忆长期以守城与屠城叙述该事件} 留下本节点追踪「江阴在薙发命令和清军围攻中失守，地方记忆长期以守城与屠城叙述该事件」之后可见的余波的材料坐标。清军与江阴地方守军所面对的条件是服制命令与守备危机把行政服从转成军事对抗，清军需确保沿江交通和粮道。《清世祖实录》顺治二年江南军报；《清史稿·世祖本纪》；《江阴县志》及《江阴城守纪》支持的结果为清廷以军事压制恢复沿江节点；江阴成为地方忠义记忆，伤亡规模不能由单一叙事定量。原有置信注记：城陷主干可靠；围城天数、伤亡和人物言行多依赖地方后出文本

\paragraph{顺治二年闰六月至八月（1645年）：江阴在薙发命令和清军围攻中失守，地方记忆长期以守城与屠城叙述该事件}
\eventanchor{EQ-1645-JIANGYIN-DECISION}{江阴在薙发命令和清军围攻中失守，地方记忆长期以守城与屠城叙述该事件} 记录清军与江阴地方守军在省、府县与地方社会相互牵动的尺度的一次变化。本节点定位围绕「江阴在薙发命令和清军围攻中失守，地方记忆长期以守城与屠城叙述该事件」形成的处置与调度记录。其发生与服制命令与守备危机把行政服从转成军事对抗，清军需确保沿江交通和粮道相连；《清世祖实录》顺治二年江南军报；《清史稿·世祖本纪》；《江阴县志》及《江阴城守纪》把事件系回来源，清廷以军事压制恢复沿江节点；江阴成为地方忠义记忆，伤亡规模不能由单一叙事定量显示压力如何向后转移。原有置信注记：城陷主干可靠；围城天数、伤亡和人物言行多依赖地方后出文本

\paragraph{顺治二年闰六月至八月（1645年）：江阴在薙发命令和清军围攻中失守，地方记忆长期以守城与屠城叙述该事件}
\eventanchor{EQ-1645-JIANGYIN-PRELUDE}{江阴在薙发命令和清军围攻中失守，地方记忆长期以守城与屠城叙述该事件} 的地点与行动范围属于省、府县与地方社会相互牵动的尺度，行动者是清军与江阴地方守军。本节点回看「江阴在薙发命令和清军围攻中失守，地方记忆长期以守城与屠城叙述该事件」之前已可见的条件。服制命令与守备危机把行政服从转成军事对抗，清军需确保沿江交通和粮道。读《清世祖实录》顺治二年江南军报；《清史稿·世祖本纪》；《江阴县志》及《江阴城守纪》时可见其记录边界，最小后果只能写作清廷以军事压制恢复沿江节点；江阴成为地方忠义记忆，伤亡规模不能由单一叙事定量。原有置信注记：城陷主干可靠；围城天数、伤亡和人物言行多依赖地方后出文本

\paragraph{顺治二年四月（1645年5月前后）：多铎军攻取南京，南明弘光朝的都城失守}
\eventanchor{EQ-1645-NANJING-TAKEN}{多铎军攻取南京，南明弘光朝的都城失守} 是清与南明弘光政权在京师与多地军政线路相接的尺度留下的编年标记。本节点叙述该事件本身。它从清军沿运河与淮扬推进，弘光朝党争、军饷困难和江防失守相互叠加走来；《清世祖实录》顺治二年四月；《清史稿·世祖本纪》；《南明史料》所收弘光朝奏疏与江宁地方记事固定可证部分，南明失去首都和财政文化中心，江浙、闽粤的抵抗及另立政权随之分化构成下一步的可见结果。原有置信注记：城陷日期可靠；守军数量、降附动机和战后暴力须并读南明与地方材料

\paragraph{顺治二年四月（1645年5月前后）：多铎军攻取南京，南明弘光朝的都城失守}
本节点追踪「多铎军攻取南京，南明弘光朝的都城失守」之后可见的余波：\eventanchor{EQ-1645-NANJING-TAKEN-AFTERMATH}{多铎军攻取南京，南明弘光朝的都城失守}。清与南明弘光政权在京师与多地军政线路相接的尺度的处境可由清军沿运河与淮扬推进，弘光朝党争、军饷困难和江防失守相互叠加说明。《清世祖实录》顺治二年四月；《清史稿·世祖本纪》；《南明史料》所收弘光朝奏疏与江宁地方记事是本段材料入口，因而只能据此追踪到南明失去首都和财政文化中心，江浙、闽粤的抵抗及另立政权随之分化。原有置信注记：城陷日期可靠；守军数量、降附动机和战后暴力须并读南明与地方材料

\paragraph{顺治二年四月（1645年5月前后）：多铎军攻取南京，南明弘光朝的都城失守}
\eventanchor{EQ-1645-NANJING-TAKEN-DECISION}{多铎军攻取南京，南明弘光朝的都城失守} 发生在京师与多地军政线路相接的尺度，牵动清与南明弘光政权。本节点定位围绕「多铎军攻取南京，南明弘光朝的都城失守」形成的处置与调度记录。清军沿运河与淮扬推进，弘光朝党争、军饷困难和江防失守相互叠加是其结构性条件；《清世祖实录》顺治二年四月；《清史稿·世祖本纪》；《南明史料》所收弘光朝奏疏与江宁地方记事留下的证据把读者带到南明失去首都和财政文化中心，江浙、闽粤的抵抗及另立政权随之分化。原有置信注记：城陷日期可靠；守军数量、降附动机和战后暴力须并读南明与地方材料

\paragraph{顺治二年四月（1645年5月前后）：多铎军攻取南京，南明弘光朝的都城失守}
\eventanchor{EQ-1645-NANJING-TAKEN-PRELUDE}{多铎军攻取南京，南明弘光朝的都城失守} 对应本节点回看「多铎军攻取南京，南明弘光朝的都城失守」之前已可见的条件。在京师与多地军政线路相接的尺度，清与南明弘光政权无法绕开清军沿运河与淮扬推进，弘光朝党争、军饷困难和江防失守相互叠加。依据《清世祖实录》顺治二年四月；《清史稿·世祖本纪》；《南明史料》所收弘光朝奏疏与江宁地方记事，本段把后续影响限定为南明失去首都和财政文化中心，江浙、闽粤的抵抗及另立政权随之分化。原有置信注记：城陷日期可靠；守军数量、降附动机和战后暴力须并读南明与地方材料

\paragraph{顺治二年四—五月（1645年）：扬州攻陷后发生大规模杀戮与掳掠，规模和持续时间存在重大争议}
\eventanchor{EQ-1645-YANGZHOU}{扬州攻陷后发生大规模杀戮与掳掠，规模和持续时间存在重大争议}。本节点叙述该事件本身。清军与扬州在京师与多地军政线路相接的尺度的处境首先受扬州守军扼守清军南下通道，攻城后的军纪、报复和掠夺处于战时权力失控与奖惩机制中制约。《清世祖实录》顺治二年四月扬州军报；《清史稿·史可法传》；王秀楚《扬州十日记》与《扬州府志》留下可核的线索；可辨的后果是扬州经验在江南传播，既促成恐惧性降附，也强化若干城镇抵抗及后世纪念政治。原有置信注记：扬州失陷可靠；死亡数字和“十日”叙事须标为后出见证并与地方材料互证

\paragraph{顺治二年四—五月（1645年）：扬州攻陷后发生大规模杀戮与掳掠，规模和持续时间存在重大争议}
\eventanchor{EQ-1645-YANGZHOU-AFTERMATH}{扬州攻陷后发生大规模杀戮与掳掠，规模和持续时间存在重大争议} 置于京师与多地军政线路相接的尺度中阅读，本节点追踪「扬州攻陷后发生大规模杀戮与掳掠，规模和持续时间存在重大争议」之后可见的余波。扬州守军扼守清军南下通道，攻城后的军纪、报复和掠夺处于战时权力失控与奖惩机制中构成行动者清军与扬州必须面对的条件。取证从《清世祖实录》顺治二年四月扬州军报；《清史稿·史可法传》；王秀楚《扬州十日记》与《扬州府志》出发，因而只能把变化谨慎地连到扬州经验在江南传播，既促成恐惧性降附，也强化若干城镇抵抗及后世纪念政治。原有置信注记：扬州失陷可靠；死亡数字和“十日”叙事须标为后出见证并与地方材料互证

\paragraph{顺治二年四—五月（1645年）：扬州攻陷后发生大规模杀戮与掳掠，规模和持续时间存在重大争议}
\eventanchor{EQ-1645-YANGZHOU-DECISION}{扬州攻陷后发生大规模杀戮与掳掠，规模和持续时间存在重大争议} 标记本节点定位围绕「扬州攻陷后发生大规模杀戮与掳掠，规模和持续时间存在重大争议」形成的处置与调度记录，涉及清军与扬州。京师与多地军政线路相接的尺度不是抽象背景：扬州守军扼守清军南下通道，攻城后的军纪、报复和掠夺处于战时权力失控与奖惩机制中。《清世祖实录》顺治二年四月扬州军报；《清史稿·史可法传》；王秀楚《扬州十日记》与《扬州府志》所见的顺序随后指向扬州经验在江南传播，既促成恐惧性降附，也强化若干城镇抵抗及后世纪念政治，但原有置信注记：扬州失陷可靠；死亡数字和“十日”叙事须标为后出见证并与地方材料互证

\paragraph{顺治二年四—五月（1645年）：扬州攻陷后发生大规模杀戮与掳掠，规模和持续时间存在重大争议}
在京师与多地军政线路相接的尺度，\eventanchor{EQ-1645-YANGZHOU-PRELUDE}{扬州攻陷后发生大规模杀戮与掳掠，规模和持续时间存在重大争议} 让清军与扬州的一个选择或压力有了可查的坐标。本节点回看「扬州攻陷后发生大规模杀戮与掳掠，规模和持续时间存在重大争议」之前已可见的条件。前因可写为扬州守军扼守清军南下通道，攻城后的军纪、报复和掠夺处于战时权力失控与奖惩机制中；《清世祖实录》顺治二年四月扬州军报；《清史稿·史可法传》；王秀楚《扬州十日记》与《扬州府志》限定了记录，后续变化是扬州经验在江南传播，既促成恐惧性降附，也强化若干城镇抵抗及后世纪念政治。原有置信注记：扬州失陷可靠；死亡数字和“十日”叙事须标为后出见证并与地方材料互证

\paragraph{顺治二年六月至十月（1645年）：清军进取杭州、绍兴等地，南明鲁王政权转入浙东及海上活动}
\eventanchor{EQ-1645-ZHEJIANG-CONQUEST}{清军进取杭州、绍兴等地，南明鲁王政权转入浙东及海上活动} 所关联的是清与南明鲁王集团、浙江地方势力在港口、海峡与跨海航路相连的尺度中的行动。本节点叙述该事件本身。它从南京失守后浙东军民与明宗室试图保持独立指挥，清军沿水陆节点分割其资源展开，借《清世祖实录》顺治二年六月至十月；《清史稿·鲁王传》；《浙江通志》与鲁王政权文书辑录进入可检验的年序，并以清廷取得浙江大部城镇，鲁王集团与郑氏网络相连，海疆战争成为长期问题影响下一段安排。原有置信注记：城市失守与鲁王转移可靠；各府实际控制和海岛联系须逐地核查

\paragraph{顺治二年六月至十月（1645年）：清军进取杭州、绍兴等地，南明鲁王政权转入浙东及海上活动}
港口、海峡与跨海航路相连的尺度中的\eventanchor{EQ-1645-ZHEJIANG-CONQUEST-AFTERMATH}{清军进取杭州、绍兴等地，南明鲁王政权转入浙东及海上活动} 不能离开清与南明鲁王集团、浙江地方势力来解释。本节点追踪「清军进取杭州、绍兴等地，南明鲁王政权转入浙东及海上活动」之后可见的余波。南京失守后浙东军民与明宗室试图保持独立指挥，清军沿水陆节点分割其资源说明其形成的条件；《清世祖实录》顺治二年六月至十月；《清史稿·鲁王传》；《浙江通志》与鲁王政权文书辑录提供来源路径，清廷取得浙江大部城镇，鲁王集团与郑氏网络相连，海疆战争成为长期问题则是材料能够支持的近时结果。原有置信注记：城市失守与鲁王转移可靠；各府实际控制和海岛联系须逐地核查

\paragraph{顺治二年六月至十月（1645年）：清军进取杭州、绍兴等地，南明鲁王政权转入浙东及海上活动}
\eventanchor{EQ-1645-ZHEJIANG-CONQUEST-DECISION}{清军进取杭州、绍兴等地，南明鲁王政权转入浙东及海上活动} 把清与南明鲁王集团、浙江地方势力放回港口、海峡与跨海航路相连的尺度的道路、城镇、边界或制度网络。本节点定位围绕「清军进取杭州、绍兴等地，南明鲁王政权转入浙东及海上活动」形成的处置与调度记录。南京失守后浙东军民与明宗室试图保持独立指挥，清军沿水陆节点分割其资源。本段采用《清世祖实录》顺治二年六月至十月；《清史稿·鲁王传》；《浙江通志》与鲁王政权文书辑录核对其顺序，随后可追到清廷取得浙江大部城镇，鲁王集团与郑氏网络相连，海疆战争成为长期问题。原有置信注记：城市失守与鲁王转移可靠；各府实际控制和海岛联系须逐地核查

\paragraph{顺治二年六月至十月（1645年）：清军进取杭州、绍兴等地，南明鲁王政权转入浙东及海上活动}
本节点回看「清军进取杭州、绍兴等地，南明鲁王政权转入浙东及海上活动」之前已可见的条件，故\eventanchor{EQ-1645-ZHEJIANG-CONQUEST-PRELUDE}{清军进取杭州、绍兴等地，南明鲁王政权转入浙东及海上活动} 不只是一个年份标签。清与南明鲁王集团、浙江地方势力在港口、海峡与跨海航路相连的尺度承受的压力来自南京失守后浙东军民与明宗室试图保持独立指挥，清军沿水陆节点分割其资源。《清世祖实录》顺治二年六月至十月；《清史稿·鲁王传》；《浙江通志》与鲁王政权文书辑录可回查该节点；清廷取得浙江大部城镇，鲁王集团与郑氏网络相连，海疆战争成为长期问题是其进入后续年序的方式。原有置信注记：城市失守与鲁王转移可靠；各府实际控制和海岛联系须逐地核查

\paragraph{顺治三年八月（1646年）：清军在闽北俘杀隆武帝朱聿键，隆武朝终结}
\eventanchor{EQ-1646-LONGWU-CAPTURE}{清军在闽北俘杀隆武帝朱聿键，隆武朝终结} 所见的行动者为清与南明隆武政权，空间尺度为京师与多地军政线路相接的尺度。本节点叙述该事件本身。清军由浙赣入闽，隆武政权军饷和将领协调不足且难以掌握山海通道使它成为具体事件链的一环；《清世祖实录》顺治三年八月；《清史稿·隆武帝本纪》；《福建通志》与黄道周文集相关记载与闽中明廷名义被打断，郑氏网络和地方武装成为东南抗清的重要行动者。原有置信注记：被俘死亡主干可靠；行止路线与地方助力记载不一

\paragraph{顺治三年八月（1646年）：清军在闽北俘杀隆武帝朱聿键，隆武朝终结}
从京师与多地军政线路相接的尺度看，\eventanchor{EQ-1646-LONGWU-CAPTURE-AFTERMATH}{清军在闽北俘杀隆武帝朱聿键，隆武朝终结} 留下本节点追踪「清军在闽北俘杀隆武帝朱聿键，隆武朝终结」之后可见的余波的材料坐标。清与南明隆武政权所面对的条件是清军由浙赣入闽，隆武政权军饷和将领协调不足且难以掌握山海通道。《清世祖实录》顺治三年八月；《清史稿·隆武帝本纪》；《福建通志》与黄道周文集相关记载支持的结果为闽中明廷名义被打断，郑氏网络和地方武装成为东南抗清的重要行动者。原有置信注记：被俘死亡主干可靠；行止路线与地方助力记载不一

\paragraph{顺治三年八月（1646年）：清军在闽北俘杀隆武帝朱聿键，隆武朝终结}
\eventanchor{EQ-1646-LONGWU-CAPTURE-DECISION}{清军在闽北俘杀隆武帝朱聿键，隆武朝终结} 记录清与南明隆武政权在京师与多地军政线路相接的尺度的一次变化。本节点定位围绕「清军在闽北俘杀隆武帝朱聿键，隆武朝终结」形成的处置与调度记录。其发生与清军由浙赣入闽，隆武政权军饷和将领协调不足且难以掌握山海通道相连；《清世祖实录》顺治三年八月；《清史稿·隆武帝本纪》；《福建通志》与黄道周文集相关记载把事件系回来源，闽中明廷名义被打断，郑氏网络和地方武装成为东南抗清的重要行动者显示压力如何向后转移。原有置信注记：被俘死亡主干可靠；行止路线与地方助力记载不一

\paragraph{顺治三年八月（1646年）：清军在闽北俘杀隆武帝朱聿键，隆武朝终结}
\eventanchor{EQ-1646-LONGWU-CAPTURE-PRELUDE}{清军在闽北俘杀隆武帝朱聿键，隆武朝终结} 的地点与行动范围属于京师与多地军政线路相接的尺度，行动者是清与南明隆武政权。本节点回看「清军在闽北俘杀隆武帝朱聿键，隆武朝终结」之前已可见的条件。清军由浙赣入闽，隆武政权军饷和将领协调不足且难以掌握山海通道。读《清世祖实录》顺治三年八月；《清史稿·隆武帝本纪》；《福建通志》与黄道周文集相关记载时可见其记录边界，最小后果只能写作闽中明廷名义被打断，郑氏网络和地方武装成为东南抗清的重要行动者。原有置信注记：被俘死亡主干可靠；行止路线与地方助力记载不一

\paragraph{顺治三年十一月（1646年）：清军攻入广州，绍武帝朱聿𨮁被杀，广州短暂的绍武朝结束}
\eventanchor{EQ-1646-SHAOWU-GUANGZHOU}{清军攻入广州，绍武帝朱聿𨮁被杀，广州短暂的绍武朝结束} 是清与南明绍武政权在京师与多地军政线路相接的尺度留下的编年标记。本节点叙述该事件本身。它从隆武朝覆灭后明宗室在广东争夺名义中心，内部冲突削弱对清军进攻的协调走来；《清世祖实录》顺治三年十一月；《清史稿·绍武帝本纪》；《广东通志》及南明文集固定可证部分，永历政权成为西南主要明廷名义，广东沿海郑氏与地方势力继续影响战局构成下一步的可见结果。原有置信注记：广州失守可靠；绍武与永历两朝冲突细节主要来自后出南明叙事

\paragraph{顺治三年十一月（1646年）：清军攻入广州，绍武帝朱聿𨮁被杀，广州短暂的绍武朝结束}
本节点追踪「清军攻入广州，绍武帝朱聿𨮁被杀，广州短暂的绍武朝结束」之后可见的余波：\eventanchor{EQ-1646-SHAOWU-GUANGZHOU-AFTERMATH}{清军攻入广州，绍武帝朱聿𨮁被杀，广州短暂的绍武朝结束}。清与南明绍武政权在京师与多地军政线路相接的尺度的处境可由隆武朝覆灭后明宗室在广东争夺名义中心，内部冲突削弱对清军进攻的协调说明。《清世祖实录》顺治三年十一月；《清史稿·绍武帝本纪》；《广东通志》及南明文集是本段材料入口，因而只能据此追踪到永历政权成为西南主要明廷名义，广东沿海郑氏与地方势力继续影响战局。原有置信注记：广州失守可靠；绍武与永历两朝冲突细节主要来自后出南明叙事

\paragraph{顺治三年十一月（1646年）：清军攻入广州，绍武帝朱聿𨮁被杀，广州短暂的绍武朝结束}
\eventanchor{EQ-1646-SHAOWU-GUANGZHOU-DECISION}{清军攻入广州，绍武帝朱聿𨮁被杀，广州短暂的绍武朝结束} 发生在京师与多地军政线路相接的尺度，牵动清与南明绍武政权。本节点定位围绕「清军攻入广州，绍武帝朱聿𨮁被杀，广州短暂的绍武朝结束」形成的处置与调度记录。隆武朝覆灭后明宗室在广东争夺名义中心，内部冲突削弱对清军进攻的协调是其结构性条件；《清世祖实录》顺治三年十一月；《清史稿·绍武帝本纪》；《广东通志》及南明文集留下的证据把读者带到永历政权成为西南主要明廷名义，广东沿海郑氏与地方势力继续影响战局。原有置信注记：广州失守可靠；绍武与永历两朝冲突细节主要来自后出南明叙事

\paragraph{顺治三年十一月（1646年）：清军攻入广州，绍武帝朱聿𨮁被杀，广州短暂的绍武朝结束}
\eventanchor{EQ-1646-SHAOWU-GUANGZHOU-PRELUDE}{清军攻入广州，绍武帝朱聿𨮁被杀，广州短暂的绍武朝结束} 对应本节点回看「清军攻入广州，绍武帝朱聿𨮁被杀，广州短暂的绍武朝结束」之前已可见的条件。在京师与多地军政线路相接的尺度，清与南明绍武政权无法绕开隆武朝覆灭后明宗室在广东争夺名义中心，内部冲突削弱对清军进攻的协调。依据《清世祖实录》顺治三年十一月；《清史稿·绍武帝本纪》；《广东通志》及南明文集，本段把后续影响限定为永历政权成为西南主要明廷名义，广东沿海郑氏与地方势力继续影响战局。原有置信注记：广州失守可靠；绍武与永历两朝冲突细节主要来自后出南明叙事

\paragraph{永历元年／顺治三年（1646年冬）：朱由榔在肇庆即位为永历帝，建立延续至西南的南明朝廷}
\eventanchor{EQ-1646-YONGLI-ZHAOQING}{朱由榔在肇庆即位为永历帝，建立延续至西南的南明朝廷}。本节点叙述该事件本身。南明永历政权与清在京师与多地军政线路相接的尺度的处境首先受绍武朝覆灭及明宗室分散，迫使支持者在两广重建朝廷；西南军政资源提供退守条件制约。《清世祖实录》顺治三年两广军报；《清史稿·永历帝本纪》；《永历实录》与《肇庆府志》留下可核的线索；可辨的后果是永历年号成为西南抗清行动的共同名义，清军须把两广、贵州、云南与缅甸路线连为战场。原有置信注记：即位时间大体可靠；改元与礼仪细节应按永历文书和清廷军报分别处理

\paragraph{永历元年／顺治三年（1646年冬）：朱由榔在肇庆即位为永历帝，建立延续至西南的南明朝廷}
\eventanchor{EQ-1646-YONGLI-ZHAOQING-AFTERMATH}{朱由榔在肇庆即位为永历帝，建立延续至西南的南明朝廷} 置于京师与多地军政线路相接的尺度中阅读，本节点追踪「朱由榔在肇庆即位为永历帝，建立延续至西南的南明朝廷」之后可见的余波。绍武朝覆灭及明宗室分散，迫使支持者在两广重建朝廷；西南军政资源提供退守条件构成行动者南明永历政权与清必须面对的条件。取证从《清世祖实录》顺治三年两广军报；《清史稿·永历帝本纪》；《永历实录》与《肇庆府志》出发，因而只能把变化谨慎地连到永历年号成为西南抗清行动的共同名义，清军须把两广、贵州、云南与缅甸路线连为战场。原有置信注记：即位时间大体可靠；改元与礼仪细节应按永历文书和清廷军报分别处理

\paragraph{永历元年／顺治三年（1646年冬）：朱由榔在肇庆即位为永历帝，建立延续至西南的南明朝廷}
\eventanchor{EQ-1646-YONGLI-ZHAOQING-DECISION}{朱由榔在肇庆即位为永历帝，建立延续至西南的南明朝廷} 标记本节点定位围绕「朱由榔在肇庆即位为永历帝，建立延续至西南的南明朝廷」形成的处置与调度记录，涉及南明永历政权与清。京师与多地军政线路相接的尺度不是抽象背景：绍武朝覆灭及明宗室分散，迫使支持者在两广重建朝廷；西南军政资源提供退守条件。《清世祖实录》顺治三年两广军报；《清史稿·永历帝本纪》；《永历实录》与《肇庆府志》所见的顺序随后指向永历年号成为西南抗清行动的共同名义，清军须把两广、贵州、云南与缅甸路线连为战场，但原有置信注记：即位时间大体可靠；改元与礼仪细节应按永历文书和清廷军报分别处理

\paragraph{永历元年／顺治三年（1646年冬）：朱由榔在肇庆即位为永历帝，建立延续至西南的南明朝廷}
在京师与多地军政线路相接的尺度，\eventanchor{EQ-1646-YONGLI-ZHAOQING-PRELUDE}{朱由榔在肇庆即位为永历帝，建立延续至西南的南明朝廷} 让南明永历政权与清的一个选择或压力有了可查的坐标。本节点回看「朱由榔在肇庆即位为永历帝，建立延续至西南的南明朝廷」之前已可见的条件。前因可写为绍武朝覆灭及明宗室分散，迫使支持者在两广重建朝廷；西南军政资源提供退守条件；《清世祖实录》顺治三年两广军报；《清史稿·永历帝本纪》；《永历实录》与《肇庆府志》限定了记录，后续变化是永历年号成为西南抗清行动的共同名义，清军须把两广、贵州、云南与缅甸路线连为战场。原有置信注记：即位时间大体可靠；改元与礼仪细节应按永历文书和清廷军报分别处理

\paragraph{顺治五年（1648年）：南昌守将金声桓反清并改奉南明，江西战局重新打开}
\eventanchor{EQ-1648-JIN-SHENGHUAN}{南昌守将金声桓反清并改奉南明，江西战局重新打开} 所关联的是清、江西与金声桓集团在京师与多地军政线路相接的尺度中的行动。本节点叙述该事件本身。它从清初任用降将以快速覆盖江南，但军饷、地位和地方网络未稳定；南明仍可提供合法性资源展开，借《清世祖实录》顺治五年江西军报；《清史稿·金声桓传》；《南昌府志》与南明永历朝文书进入可检验的年序，并以清廷被迫向江西投入围城和调兵资源，显示占领区行政接收不等于长期服从影响下一段安排。原有置信注记：反叛及南昌控制可靠；部众规模和改奉动机多为交战各方归责语言

\paragraph{顺治五年（1648年）：南昌守将金声桓反清并改奉南明，江西战局重新打开}
京师与多地军政线路相接的尺度中的\eventanchor{EQ-1648-JIN-SHENGHUAN-AFTERMATH}{南昌守将金声桓反清并改奉南明，江西战局重新打开} 不能离开清、江西与金声桓集团来解释。本节点追踪「南昌守将金声桓反清并改奉南明，江西战局重新打开」之后可见的余波。清初任用降将以快速覆盖江南，但军饷、地位和地方网络未稳定；南明仍可提供合法性资源说明其形成的条件；《清世祖实录》顺治五年江西军报；《清史稿·金声桓传》；《南昌府志》与南明永历朝文书提供来源路径，清廷被迫向江西投入围城和调兵资源，显示占领区行政接收不等于长期服从则是材料能够支持的近时结果。原有置信注记：反叛及南昌控制可靠；部众规模和改奉动机多为交战各方归责语言

\paragraph{顺治五年（1648年）：南昌守将金声桓反清并改奉南明，江西战局重新打开}
\eventanchor{EQ-1648-JIN-SHENGHUAN-DECISION}{南昌守将金声桓反清并改奉南明，江西战局重新打开} 把清、江西与金声桓集团放回京师与多地军政线路相接的尺度的道路、城镇、边界或制度网络。本节点定位围绕「南昌守将金声桓反清并改奉南明，江西战局重新打开」形成的处置与调度记录。清初任用降将以快速覆盖江南，但军饷、地位和地方网络未稳定；南明仍可提供合法性资源。本段采用《清世祖实录》顺治五年江西军报；《清史稿·金声桓传》；《南昌府志》与南明永历朝文书核对其顺序，随后可追到清廷被迫向江西投入围城和调兵资源，显示占领区行政接收不等于长期服从。原有置信注记：反叛及南昌控制可靠；部众规模和改奉动机多为交战各方归责语言

\paragraph{顺治五年（1648年）：南昌守将金声桓反清并改奉南明，江西战局重新打开}
本节点回看「南昌守将金声桓反清并改奉南明，江西战局重新打开」之前已可见的条件，故\eventanchor{EQ-1648-JIN-SHENGHUAN-PRELUDE}{南昌守将金声桓反清并改奉南明，江西战局重新打开} 不只是一个年份标签。清、江西与金声桓集团在京师与多地军政线路相接的尺度承受的压力来自清初任用降将以快速覆盖江南，但军饷、地位和地方网络未稳定；南明仍可提供合法性资源。《清世祖实录》顺治五年江西军报；《清史稿·金声桓传》；《南昌府志》与南明永历朝文书可回查该节点；清廷被迫向江西投入围城和调兵资源，显示占领区行政接收不等于长期服从是其进入后续年序的方式。原有置信注记：反叛及南昌控制可靠；部众规模和改奉动机多为交战各方归责语言

\paragraph{顺治五年（1648年）：广东清将李成栋反清，广州一度转入永历朝控制}
\eventanchor{EQ-1648-LI-CHENGDONG}{广东清将李成栋反清，广州一度转入永历朝控制} 所见的行动者为清、广东与李成栋集团，空间尺度为省、府县与地方社会相互牵动的尺度。本节点叙述该事件本身。清军依靠原明将领维系两广，军功分配与中央控制的矛盾给永历朝重新联络的空间使它成为具体事件链的一环；《清世祖实录》顺治五年广东军报；《清史稿·李成栋传》；《广东通志》及永历朝诏敕辑录与两广战场扩大并与江西、闽海相连，清廷须重新部署南方军饷和将领体系。原有置信注记：反叛与广州易手可靠；个人动机和地方支持范围有强烈阵营叙事差异

\paragraph{顺治五年（1648年）：广东清将李成栋反清，广州一度转入永历朝控制}
从省、府县与地方社会相互牵动的尺度看，\eventanchor{EQ-1648-LI-CHENGDONG-AFTERMATH}{广东清将李成栋反清，广州一度转入永历朝控制} 留下本节点追踪「广东清将李成栋反清，广州一度转入永历朝控制」之后可见的余波的材料坐标。清、广东与李成栋集团所面对的条件是清军依靠原明将领维系两广，军功分配与中央控制的矛盾给永历朝重新联络的空间。《清世祖实录》顺治五年广东军报；《清史稿·李成栋传》；《广东通志》及永历朝诏敕辑录支持的结果为两广战场扩大并与江西、闽海相连，清廷须重新部署南方军饷和将领体系。原有置信注记：反叛与广州易手可靠；个人动机和地方支持范围有强烈阵营叙事差异

\paragraph{顺治五年（1648年）：广东清将李成栋反清，广州一度转入永历朝控制}
\eventanchor{EQ-1648-LI-CHENGDONG-DECISION}{广东清将李成栋反清，广州一度转入永历朝控制} 记录清、广东与李成栋集团在省、府县与地方社会相互牵动的尺度的一次变化。本节点定位围绕「广东清将李成栋反清，广州一度转入永历朝控制」形成的处置与调度记录。其发生与清军依靠原明将领维系两广，军功分配与中央控制的矛盾给永历朝重新联络的空间相连；《清世祖实录》顺治五年广东军报；《清史稿·李成栋传》；《广东通志》及永历朝诏敕辑录把事件系回来源，两广战场扩大并与江西、闽海相连，清廷须重新部署南方军饷和将领体系显示压力如何向后转移。原有置信注记：反叛与广州易手可靠；个人动机和地方支持范围有强烈阵营叙事差异

\paragraph{顺治五年（1648年）：广东清将李成栋反清，广州一度转入永历朝控制}
\eventanchor{EQ-1648-LI-CHENGDONG-PRELUDE}{广东清将李成栋反清，广州一度转入永历朝控制} 的地点与行动范围属于省、府县与地方社会相互牵动的尺度，行动者是清、广东与李成栋集团。本节点回看「广东清将李成栋反清，广州一度转入永历朝控制」之前已可见的条件。清军依靠原明将领维系两广，军功分配与中央控制的矛盾给永历朝重新联络的空间。读《清世祖实录》顺治五年广东军报；《清史稿·李成栋传》；《广东通志》及永历朝诏敕辑录时可见其记录边界，最小后果只能写作两广战场扩大并与江西、闽海相连，清廷须重新部署南方军饷和将领体系。原有置信注记：反叛与广州易手可靠；个人动机和地方支持范围有强烈阵营叙事差异

\paragraph{顺治六年正月至三月（1649年）：清军围攻并攻取南昌，金声桓反叛失败}
\eventanchor{EQ-1649-NANCHANG}{清军围攻并攻取南昌，金声桓反叛失败} 是清与金声桓集团在京师与多地军政线路相接的尺度留下的编年标记。本节点叙述该事件本身。它从金声桓据南昌切断清廷赣江通道，清军需以持续军需和外围据点完成围攻走来；《清世祖实录》顺治六年正月至三月；《清史稿·金声桓传》；内阁大库顺治六年江西军需题本与《南昌府志》固定可证部分，江西反叛中心被击破，清廷恢复通道但须面对战后税源、流民和秩序重建构成下一步的可见结果。原有置信注记：围城结局可靠；攻城损失、城内处置和战后人口变化须以地方资料复核

\paragraph{顺治六年正月至三月（1649年）：清军围攻并攻取南昌，金声桓反叛失败}
本节点追踪「清军围攻并攻取南昌，金声桓反叛失败」之后可见的余波：\eventanchor{EQ-1649-NANCHANG-AFTERMATH}{清军围攻并攻取南昌，金声桓反叛失败}。清与金声桓集团在京师与多地军政线路相接的尺度的处境可由金声桓据南昌切断清廷赣江通道，清军需以持续军需和外围据点完成围攻说明。《清世祖实录》顺治六年正月至三月；《清史稿·金声桓传》；内阁大库顺治六年江西军需题本与《南昌府志》是本段材料入口，因而只能据此追踪到江西反叛中心被击破，清廷恢复通道但须面对战后税源、流民和秩序重建。原有置信注记：围城结局可靠；攻城损失、城内处置和战后人口变化须以地方资料复核

\paragraph{顺治六年正月至三月（1649年）：清军围攻并攻取南昌，金声桓反叛失败}
\eventanchor{EQ-1649-NANCHANG-DECISION}{清军围攻并攻取南昌，金声桓反叛失败} 发生在京师与多地军政线路相接的尺度，牵动清与金声桓集团。本节点定位围绕「清军围攻并攻取南昌，金声桓反叛失败」形成的处置与调度记录。金声桓据南昌切断清廷赣江通道，清军需以持续军需和外围据点完成围攻是其结构性条件；《清世祖实录》顺治六年正月至三月；《清史稿·金声桓传》；内阁大库顺治六年江西军需题本与《南昌府志》留下的证据把读者带到江西反叛中心被击破，清廷恢复通道但须面对战后税源、流民和秩序重建。原有置信注记：围城结局可靠；攻城损失、城内处置和战后人口变化须以地方资料复核

\paragraph{顺治六年正月至三月（1649年）：清军围攻并攻取南昌，金声桓反叛失败}
\eventanchor{EQ-1649-NANCHANG-PRELUDE}{清军围攻并攻取南昌，金声桓反叛失败} 对应本节点回看「清军围攻并攻取南昌，金声桓反叛失败」之前已可见的条件。在京师与多地军政线路相接的尺度，清与金声桓集团无法绕开金声桓据南昌切断清廷赣江通道，清军需以持续军需和外围据点完成围攻。依据《清世祖实录》顺治六年正月至三月；《清史稿·金声桓传》；内阁大库顺治六年江西军需题本与《南昌府志》，本段把后续影响限定为江西反叛中心被击破，清廷恢复通道但须面对战后税源、流民和秩序重建。原有置信注记：围城结局可靠；攻城损失、城内处置和战后人口变化须以地方资料复核

\paragraph{顺治七年十一月（1650年）：清军重新攻取广州，战后杀戮与损毁规模成为广东地方记忆中的重大争议}
\eventanchor{EQ-1650-GUANGZHOU}{清军重新攻取广州，战后杀戮与损毁规模成为广东地方记忆中的重大争议}。本节点叙述该事件本身。清与永历朝广东守军在省、府县与地方社会相互牵动的尺度的处境首先受李成栋反叛后广东成为永历朝资源节点，清军重夺广州兼具军事和海防目的制约。《清世祖实录》顺治七年十一月；《清史稿·世祖本纪》；《广东通志》、广州地方笔记与内阁大库广东军报留下可核的线索；可辨的后果是永历朝在两广基础削弱，广州损失与人口流动影响地方财政、宗族和港口恢复。原有置信注记：广州再陷可靠；死亡数字和“庚寅之劫”持续时间不能从清军战报单定

\paragraph{顺治七年十一月（1650年）：清军重新攻取广州，战后杀戮与损毁规模成为广东地方记忆中的重大争议}
\eventanchor{EQ-1650-GUANGZHOU-AFTERMATH}{清军重新攻取广州，战后杀戮与损毁规模成为广东地方记忆中的重大争议} 置于省、府县与地方社会相互牵动的尺度中阅读，本节点追踪「清军重新攻取广州，战后杀戮与损毁规模成为广东地方记忆中的重大争议」之后可见的余波。李成栋反叛后广东成为永历朝资源节点，清军重夺广州兼具军事和海防目的构成行动者清与永历朝广东守军必须面对的条件。取证从《清世祖实录》顺治七年十一月；《清史稿·世祖本纪》；《广东通志》、广州地方笔记与内阁大库广东军报出发，因而只能把变化谨慎地连到永历朝在两广基础削弱，广州损失与人口流动影响地方财政、宗族和港口恢复。原有置信注记：广州再陷可靠；死亡数字和“庚寅之劫”持续时间不能从清军战报单定

\paragraph{顺治七年十一月（1650年）：清军重新攻取广州，战后杀戮与损毁规模成为广东地方记忆中的重大争议}
\eventanchor{EQ-1650-GUANGZHOU-DECISION}{清军重新攻取广州，战后杀戮与损毁规模成为广东地方记忆中的重大争议} 标记本节点定位围绕「清军重新攻取广州，战后杀戮与损毁规模成为广东地方记忆中的重大争议」形成的处置与调度记录，涉及清与永历朝广东守军。省、府县与地方社会相互牵动的尺度不是抽象背景：李成栋反叛后广东成为永历朝资源节点，清军重夺广州兼具军事和海防目的。《清世祖实录》顺治七年十一月；《清史稿·世祖本纪》；《广东通志》、广州地方笔记与内阁大库广东军报所见的顺序随后指向永历朝在两广基础削弱，广州损失与人口流动影响地方财政、宗族和港口恢复，但原有置信注记：广州再陷可靠；死亡数字和“庚寅之劫”持续时间不能从清军战报单定

\paragraph{顺治七年十一月（1650年）：清军重新攻取广州，战后杀戮与损毁规模成为广东地方记忆中的重大争议}
在省、府县与地方社会相互牵动的尺度，\eventanchor{EQ-1650-GUANGZHOU-PRELUDE}{清军重新攻取广州，战后杀戮与损毁规模成为广东地方记忆中的重大争议} 让清与永历朝广东守军的一个选择或压力有了可查的坐标。本节点回看「清军重新攻取广州，战后杀戮与损毁规模成为广东地方记忆中的重大争议」之前已可见的条件。前因可写为李成栋反叛后广东成为永历朝资源节点，清军重夺广州兼具军事和海防目的；《清世祖实录》顺治七年十一月；《清史稿·世祖本纪》；《广东通志》、广州地方笔记与内阁大库广东军报限定了记录，后续变化是永历朝在两广基础削弱，广州损失与人口流动影响地方财政、宗族和港口恢复。原有置信注记：广州再陷可靠；死亡数字和“庚寅之劫”持续时间不能从清军战报单定

\paragraph{永历六年／顺治九年（1652年）：李定国攻克桂林，孔有德自尽，清在广西的军事据点受重创}
\eventanchor{EQ-1652-DINGGUO-GUILIN}{李定国攻克桂林，孔有德自尽，清在广西的军事据点受重创} 所关联的是南明永历朝、清与李定国军在京师与多地军政线路相接的尺度中的行动。本节点叙述该事件本身。它从两广清军战线长、补给困难且将领协调不足，永历军利用机动兵力突击桂林展开，借《清世祖实录》顺治九年广西军报；《清史稿·孔有德传》；《永历实录》与《广西通志》进入可检验的年序，并以永历朝获得声望与物资，但未能稳定保持两广；清廷随即加强西南调兵影响下一段安排。原有置信注记：桂林战果可靠；双方兵数和孔有德死因细节应保留不同文本说法

\paragraph{永历六年／顺治九年（1652年）：李定国攻克桂林，孔有德自尽，清在广西的军事据点受重创}
京师与多地军政线路相接的尺度中的\eventanchor{EQ-1652-DINGGUO-GUILIN-AFTERMATH}{李定国攻克桂林，孔有德自尽，清在广西的军事据点受重创} 不能离开南明永历朝、清与李定国军来解释。本节点追踪「李定国攻克桂林，孔有德自尽，清在广西的军事据点受重创」之后可见的余波。两广清军战线长、补给困难且将领协调不足，永历军利用机动兵力突击桂林说明其形成的条件；《清世祖实录》顺治九年广西军报；《清史稿·孔有德传》；《永历实录》与《广西通志》提供来源路径，永历朝获得声望与物资，但未能稳定保持两广；清廷随即加强西南调兵则是材料能够支持的近时结果。原有置信注记：桂林战果可靠；双方兵数和孔有德死因细节应保留不同文本说法

\paragraph{永历六年／顺治九年（1652年）：李定国攻克桂林，孔有德自尽，清在广西的军事据点受重创}
\eventanchor{EQ-1652-DINGGUO-GUILIN-DECISION}{李定国攻克桂林，孔有德自尽，清在广西的军事据点受重创} 把南明永历朝、清与李定国军放回京师与多地军政线路相接的尺度的道路、城镇、边界或制度网络。本节点定位围绕「李定国攻克桂林，孔有德自尽，清在广西的军事据点受重创」形成的处置与调度记录。两广清军战线长、补给困难且将领协调不足，永历军利用机动兵力突击桂林。本段采用《清世祖实录》顺治九年广西军报；《清史稿·孔有德传》；《永历实录》与《广西通志》核对其顺序，随后可追到永历朝获得声望与物资，但未能稳定保持两广；清廷随即加强西南调兵。原有置信注记：桂林战果可靠；双方兵数和孔有德死因细节应保留不同文本说法

\paragraph{永历六年／顺治九年（1652年）：李定国攻克桂林，孔有德自尽，清在广西的军事据点受重创}
本节点回看「李定国攻克桂林，孔有德自尽，清在广西的军事据点受重创」之前已可见的条件，故\eventanchor{EQ-1652-DINGGUO-GUILIN-PRELUDE}{李定国攻克桂林，孔有德自尽，清在广西的军事据点受重创} 不只是一个年份标签。南明永历朝、清与李定国军在京师与多地军政线路相接的尺度承受的压力来自两广清军战线长、补给困难且将领协调不足，永历军利用机动兵力突击桂林。《清世祖实录》顺治九年广西军报；《清史稿·孔有德传》；《永历实录》与《广西通志》可回查该节点；永历朝获得声望与物资，但未能稳定保持两广；清廷随即加强西南调兵是其进入后续年序的方式。原有置信注记：桂林战果可靠；双方兵数和孔有德死因细节应保留不同文本说法

\paragraph{永历六年／顺治九年夏（1652年）：李定国在衡州击败尼堪军，清军西南攻势暂受挫}
\eventanchor{EQ-1652-DINGGUO-HENGZHOU}{李定国在衡州击败尼堪军，清军西南攻势暂受挫} 所见的行动者为南明永历朝与清，空间尺度为京师与多地军政线路相接的尺度。本节点叙述该事件本身。清军试图由湖南压迫两广和贵州，李定国依托本地军队与道路选择反击使它成为具体事件链的一环；《清世祖实录》顺治九年湖南军报；《清史稿·尼堪传》；《永历实录》与《衡州府志》与清廷认识到西南征服不能只靠快速突进，转向更大规模调兵、饷运和分区推进。原有置信注记：战役结局大体可靠；地点、伤亡与尼堪死亡过程在中清叙事与南明记录间有异

\paragraph{永历六年／顺治九年夏（1652年）：李定国在衡州击败尼堪军，清军西南攻势暂受挫}
从京师与多地军政线路相接的尺度看，\eventanchor{EQ-1652-DINGGUO-HENGZHOU-AFTERMATH}{李定国在衡州击败尼堪军，清军西南攻势暂受挫} 留下本节点追踪「李定国在衡州击败尼堪军，清军西南攻势暂受挫」之后可见的余波的材料坐标。南明永历朝与清所面对的条件是清军试图由湖南压迫两广和贵州，李定国依托本地军队与道路选择反击。《清世祖实录》顺治九年湖南军报；《清史稿·尼堪传》；《永历实录》与《衡州府志》支持的结果为清廷认识到西南征服不能只靠快速突进，转向更大规模调兵、饷运和分区推进。原有置信注记：战役结局大体可靠；地点、伤亡与尼堪死亡过程在中清叙事与南明记录间有异

\paragraph{永历六年／顺治九年夏（1652年）：李定国在衡州击败尼堪军，清军西南攻势暂受挫}
\eventanchor{EQ-1652-DINGGUO-HENGZHOU-DECISION}{李定国在衡州击败尼堪军，清军西南攻势暂受挫} 记录南明永历朝与清在京师与多地军政线路相接的尺度的一次变化。本节点定位围绕「李定国在衡州击败尼堪军，清军西南攻势暂受挫」形成的处置与调度记录。其发生与清军试图由湖南压迫两广和贵州，李定国依托本地军队与道路选择反击相连；《清世祖实录》顺治九年湖南军报；《清史稿·尼堪传》；《永历实录》与《衡州府志》把事件系回来源，清廷认识到西南征服不能只靠快速突进，转向更大规模调兵、饷运和分区推进显示压力如何向后转移。原有置信注记：战役结局大体可靠；地点、伤亡与尼堪死亡过程在中清叙事与南明记录间有异

\paragraph{永历六年／顺治九年夏（1652年）：李定国在衡州击败尼堪军，清军西南攻势暂受挫}
\eventanchor{EQ-1652-DINGGUO-HENGZHOU-PRELUDE}{李定国在衡州击败尼堪军，清军西南攻势暂受挫} 的地点与行动范围属于京师与多地军政线路相接的尺度，行动者是南明永历朝与清。本节点回看「李定国在衡州击败尼堪军，清军西南攻势暂受挫」之前已可见的条件。清军试图由湖南压迫两广和贵州，李定国依托本地军队与道路选择反击。读《清世祖实录》顺治九年湖南军报；《清史稿·尼堪传》；《永历实录》与《衡州府志》时可见其记录边界，最小后果只能写作清廷认识到西南征服不能只靠快速突进，转向更大规模调兵、饷运和分区推进。原有置信注记：战役结局大体可靠；地点、伤亡与尼堪死亡过程在中清叙事与南明记录间有异

\paragraph{永历十一年／顺治十四年（1657年）：永历朝在贵州、云南的孙可望与李定国集团矛盾公开化，孙可望转降清廷}
\eventanchor{EQ-1657-YONGLI-GUIZHOU}{永历朝在贵州、云南的孙可望与李定国集团矛盾公开化，孙可望转降清廷} 是南明永历朝、清与孙可望集团在京师与多地军政线路相接的尺度留下的编年标记。本节点叙述该事件本身。它从长期战争中的饷源、官爵和指挥权竞争削弱永历朝合作，清廷提供孙可望转换阵营的条件走来；《清世祖实录》顺治十四年贵州军报；《清史稿·孙可望传》；《永历实录》与贵州地方志固定可证部分，清军获得西南军情与向导资源，永历朝失去大批可动员兵力，云南退路日益孤立构成下一步的可见结果。原有置信注记：孙可望降清可靠；冲突责任、军队规模与永历帝处境受各方政治文本影响

\paragraph{永历十一年／顺治十四年（1657年）：永历朝在贵州、云南的孙可望与李定国集团矛盾公开化，孙可望转降清廷}
本节点追踪「永历朝在贵州、云南的孙可望与李定国集团矛盾公开化，孙可望转降清廷」之后可见的余波：\eventanchor{EQ-1657-YONGLI-GUIZHOU-AFTERMATH}{永历朝在贵州、云南的孙可望与李定国集团矛盾公开化，孙可望转降清廷}。南明永历朝、清与孙可望集团在京师与多地军政线路相接的尺度的处境可由长期战争中的饷源、官爵和指挥权竞争削弱永历朝合作，清廷提供孙可望转换阵营的条件说明。《清世祖实录》顺治十四年贵州军报；《清史稿·孙可望传》；《永历实录》与贵州地方志是本段材料入口，因而只能据此追踪到清军获得西南军情与向导资源，永历朝失去大批可动员兵力，云南退路日益孤立。原有置信注记：孙可望降清可靠；冲突责任、军队规模与永历帝处境受各方政治文本影响

\paragraph{永历十一年／顺治十四年（1657年）：永历朝在贵州、云南的孙可望与李定国集团矛盾公开化，孙可望转降清廷}
\eventanchor{EQ-1657-YONGLI-GUIZHOU-DECISION}{永历朝在贵州、云南的孙可望与李定国集团矛盾公开化，孙可望转降清廷} 发生在京师与多地军政线路相接的尺度，牵动南明永历朝、清与孙可望集团。本节点定位围绕「永历朝在贵州、云南的孙可望与李定国集团矛盾公开化，孙可望转降清廷」形成的处置与调度记录。长期战争中的饷源、官爵和指挥权竞争削弱永历朝合作，清廷提供孙可望转换阵营的条件是其结构性条件；《清世祖实录》顺治十四年贵州军报；《清史稿·孙可望传》；《永历实录》与贵州地方志留下的证据把读者带到清军获得西南军情与向导资源，永历朝失去大批可动员兵力，云南退路日益孤立。原有置信注记：孙可望降清可靠；冲突责任、军队规模与永历帝处境受各方政治文本影响

\paragraph{永历十一年／顺治十四年（1657年）：永历朝在贵州、云南的孙可望与李定国集团矛盾公开化，孙可望转降清廷}
\eventanchor{EQ-1657-YONGLI-GUIZHOU-PRELUDE}{永历朝在贵州、云南的孙可望与李定国集团矛盾公开化，孙可望转降清廷} 对应本节点回看「永历朝在贵州、云南的孙可望与李定国集团矛盾公开化，孙可望转降清廷」之前已可见的条件。在京师与多地军政线路相接的尺度，南明永历朝、清与孙可望集团无法绕开长期战争中的饷源、官爵和指挥权竞争削弱永历朝合作，清廷提供孙可望转换阵营的条件。依据《清世祖实录》顺治十四年贵州军报；《清史稿·孙可望传》；《永历实录》与贵州地方志，本段把后续影响限定为清军获得西南军情与向导资源，永历朝失去大批可动员兵力，云南退路日益孤立。原有置信注记：孙可望降清可靠；冲突责任、军队规模与永历帝处境受各方政治文本影响

\paragraph{顺治十五年（1658年）：清军由贵州、四川等方向进逼云南，永历朝廷撤离昆明并向西南退却}
\eventanchor{EQ-1658-QING-YUNNAN}{清军由贵州、四川等方向进逼云南，永历朝廷撤离昆明并向西南退却}。本节点叙述该事件本身。清与南明永历朝在省、府县与地方社会相互牵动的尺度的处境首先受孙可望降清削弱南明防线，清军依托多路会攻和地方土司、降将信息推进制约。《清世祖实录》顺治十五年云南军报；《清史稿·洪承畴传》；《云南通志》与永历朝奏疏辑录留下可核的线索；可辨的后果是永历政权转为流动朝廷，清军虽占昆明等城，仍须在边地和交通线上持续作战。原有置信注记：清军入滇和永历撤离可靠；各路行军与地方归附应分区而非视作全省一次完成

\paragraph{顺治十五年（1658年）：清军由贵州、四川等方向进逼云南，永历朝廷撤离昆明并向西南退却}
\eventanchor{EQ-1658-QING-YUNNAN-AFTERMATH}{清军由贵州、四川等方向进逼云南，永历朝廷撤离昆明并向西南退却} 置于省、府县与地方社会相互牵动的尺度中阅读，本节点追踪「清军由贵州、四川等方向进逼云南，永历朝廷撤离昆明并向西南退却」之后可见的余波。孙可望降清削弱南明防线，清军依托多路会攻和地方土司、降将信息推进构成行动者清与南明永历朝必须面对的条件。取证从《清世祖实录》顺治十五年云南军报；《清史稿·洪承畴传》；《云南通志》与永历朝奏疏辑录出发，因而只能把变化谨慎地连到永历政权转为流动朝廷，清军虽占昆明等城，仍须在边地和交通线上持续作战。原有置信注记：清军入滇和永历撤离可靠；各路行军与地方归附应分区而非视作全省一次完成

\paragraph{顺治十五年（1658年）：清军由贵州、四川等方向进逼云南，永历朝廷撤离昆明并向西南退却}
\eventanchor{EQ-1658-QING-YUNNAN-DECISION}{清军由贵州、四川等方向进逼云南，永历朝廷撤离昆明并向西南退却} 标记本节点定位围绕「清军由贵州、四川等方向进逼云南，永历朝廷撤离昆明并向西南退却」形成的处置与调度记录，涉及清与南明永历朝。省、府县与地方社会相互牵动的尺度不是抽象背景：孙可望降清削弱南明防线，清军依托多路会攻和地方土司、降将信息推进。《清世祖实录》顺治十五年云南军报；《清史稿·洪承畴传》；《云南通志》与永历朝奏疏辑录所见的顺序随后指向永历政权转为流动朝廷，清军虽占昆明等城，仍须在边地和交通线上持续作战，但原有置信注记：清军入滇和永历撤离可靠；各路行军与地方归附应分区而非视作全省一次完成

\paragraph{顺治十五年（1658年）：清军由贵州、四川等方向进逼云南，永历朝廷撤离昆明并向西南退却}
在省、府县与地方社会相互牵动的尺度，\eventanchor{EQ-1658-QING-YUNNAN-PRELUDE}{清军由贵州、四川等方向进逼云南，永历朝廷撤离昆明并向西南退却} 让清与南明永历朝的一个选择或压力有了可查的坐标。本节点回看「清军由贵州、四川等方向进逼云南，永历朝廷撤离昆明并向西南退却」之前已可见的条件。前因可写为孙可望降清削弱南明防线，清军依托多路会攻和地方土司、降将信息推进；《清世祖实录》顺治十五年云南军报；《清史稿·洪承畴传》；《云南通志》与永历朝奏疏辑录限定了记录，后续变化是永历政权转为流动朝廷，清军虽占昆明等城，仍须在边地和交通线上持续作战。原有置信注记：清军入滇和永历撤离可靠；各路行军与地方归附应分区而非视作全省一次完成

\paragraph{永历十三年／顺治十六年末（1659年）：永历帝一行进入缅甸，西南战争扩展为跨境避难和外交问题}
\eventanchor{EQ-1659-YONGLI-MYANMAR}{永历帝一行进入缅甸，西南战争扩展为跨境避难和外交问题} 所关联的是南明永历朝、缅甸东吁王朝与清在边地道路、驻防与多方社群交错的尺度中的行动。本节点叙述该事件本身。它从清军推进使永历朝失去云南腹地，缅甸王朝面临接纳流亡朝廷与避免清军压力的两难展开，借《清世祖实录》顺治十六年云南军报；《清史稿·永历帝本纪》；《琉璃宫史》缅文传统及欧洲旅行记译注进入可检验的年序，并以永历政权补给与外交选择受缅甸宫廷制约，清廷将追捕与边境交涉结合影响下一段安排。原有置信注记：入缅事实可靠；路线、接待条件与缅甸朝廷决策须并读汉文和缅文/欧洲记载

\paragraph{永历十三年／顺治十六年末（1659年）：永历帝一行进入缅甸，西南战争扩展为跨境避难和外交问题}
边地道路、驻防与多方社群交错的尺度中的\eventanchor{EQ-1659-YONGLI-MYANMAR-AFTERMATH}{永历帝一行进入缅甸，西南战争扩展为跨境避难和外交问题} 不能离开南明永历朝、缅甸东吁王朝与清来解释。本节点追踪「永历帝一行进入缅甸，西南战争扩展为跨境避难和外交问题」之后可见的余波。清军推进使永历朝失去云南腹地，缅甸王朝面临接纳流亡朝廷与避免清军压力的两难说明其形成的条件；《清世祖实录》顺治十六年云南军报；《清史稿·永历帝本纪》；《琉璃宫史》缅文传统及欧洲旅行记译注提供来源路径，永历政权补给与外交选择受缅甸宫廷制约，清廷将追捕与边境交涉结合则是材料能够支持的近时结果。原有置信注记：入缅事实可靠；路线、接待条件与缅甸朝廷决策须并读汉文和缅文/欧洲记载

\paragraph{永历十三年／顺治十六年末（1659年）：永历帝一行进入缅甸，西南战争扩展为跨境避难和外交问题}
\eventanchor{EQ-1659-YONGLI-MYANMAR-DECISION}{永历帝一行进入缅甸，西南战争扩展为跨境避难和外交问题} 把南明永历朝、缅甸东吁王朝与清放回边地道路、驻防与多方社群交错的尺度的道路、城镇、边界或制度网络。本节点定位围绕「永历帝一行进入缅甸，西南战争扩展为跨境避难和外交问题」形成的处置与调度记录。清军推进使永历朝失去云南腹地，缅甸王朝面临接纳流亡朝廷与避免清军压力的两难。本段采用《清世祖实录》顺治十六年云南军报；《清史稿·永历帝本纪》；《琉璃宫史》缅文传统及欧洲旅行记译注核对其顺序，随后可追到永历政权补给与外交选择受缅甸宫廷制约，清廷将追捕与边境交涉结合。原有置信注记：入缅事实可靠；路线、接待条件与缅甸朝廷决策须并读汉文和缅文/欧洲记载

\paragraph{永历十三年／顺治十六年末（1659年）：永历帝一行进入缅甸，西南战争扩展为跨境避难和外交问题}
本节点回看「永历帝一行进入缅甸，西南战争扩展为跨境避难和外交问题」之前已可见的条件，故\eventanchor{EQ-1659-YONGLI-MYANMAR-PRELUDE}{永历帝一行进入缅甸，西南战争扩展为跨境避难和外交问题} 不只是一个年份标签。南明永历朝、缅甸东吁王朝与清在边地道路、驻防与多方社群交错的尺度承受的压力来自清军推进使永历朝失去云南腹地，缅甸王朝面临接纳流亡朝廷与避免清军压力的两难。《清世祖实录》顺治十六年云南军报；《清史稿·永历帝本纪》；《琉璃宫史》缅文传统及欧洲旅行记译注可回查该节点；永历政权补给与外交选择受缅甸宫廷制约，清廷将追捕与边境交涉结合是其进入后续年序的方式。原有置信注记：入缅事实可靠；路线、接待条件与缅甸朝廷决策须并读汉文和缅文/欧洲记载

\paragraph{顺治十七年（1660年）：清军继续经营云南城镇与驿路，永历朝残部向中缅边地退缩}
\eventanchor{EQ-1660-QING-YUNNAN-CONSOLIDATION}{清军继续经营云南城镇与驿路，永历朝残部向中缅边地退缩} 所见的行动者为清与南明永历朝残部，空间尺度为京师与多地军政线路相接的尺度。本节点叙述该事件本身。占领昆明后清军仍须供应远征军、维持道路并争取地方中介，永历残部利用边地机动使它成为具体事件链的一环；《清世祖实录》顺治十七年云南军报；《清史稿·吴三桂传》；内阁大库顺治十七年云南粮饷题本与《云南通志》与云南成为吴三桂等将领驻军和财政调度基地，也为后来的藩镇结构埋下条件。原有置信注记：清军控制主要城镇大体可靠；土司、矿区和山地实际控制不能以府城易手推定

\paragraph{顺治十七年（1660年）：清军继续经营云南城镇与驿路，永历朝残部向中缅边地退缩}
从京师与多地军政线路相接的尺度看，\eventanchor{EQ-1660-QING-YUNNAN-CONSOLIDATION-AFTERMATH}{清军继续经营云南城镇与驿路，永历朝残部向中缅边地退缩} 留下本节点追踪「清军继续经营云南城镇与驿路，永历朝残部向中缅边地退缩」之后可见的余波的材料坐标。清与南明永历朝残部所面对的条件是占领昆明后清军仍须供应远征军、维持道路并争取地方中介，永历残部利用边地机动。《清世祖实录》顺治十七年云南军报；《清史稿·吴三桂传》；内阁大库顺治十七年云南粮饷题本与《云南通志》支持的结果为云南成为吴三桂等将领驻军和财政调度基地，也为后来的藩镇结构埋下条件。原有置信注记：清军控制主要城镇大体可靠；土司、矿区和山地实际控制不能以府城易手推定

\paragraph{顺治十七年（1660年）：清军继续经营云南城镇与驿路，永历朝残部向中缅边地退缩}
\eventanchor{EQ-1660-QING-YUNNAN-CONSOLIDATION-DECISION}{清军继续经营云南城镇与驿路，永历朝残部向中缅边地退缩} 记录清与南明永历朝残部在京师与多地军政线路相接的尺度的一次变化。本节点定位围绕「清军继续经营云南城镇与驿路，永历朝残部向中缅边地退缩」形成的处置与调度记录。其发生与占领昆明后清军仍须供应远征军、维持道路并争取地方中介，永历残部利用边地机动相连；《清世祖实录》顺治十七年云南军报；《清史稿·吴三桂传》；内阁大库顺治十七年云南粮饷题本与《云南通志》把事件系回来源，云南成为吴三桂等将领驻军和财政调度基地，也为后来的藩镇结构埋下条件显示压力如何向后转移。原有置信注记：清军控制主要城镇大体可靠；土司、矿区和山地实际控制不能以府城易手推定

\paragraph{顺治十七年（1660年）：清军继续经营云南城镇与驿路，永历朝残部向中缅边地退缩}
\eventanchor{EQ-1660-QING-YUNNAN-CONSOLIDATION-PRELUDE}{清军继续经营云南城镇与驿路，永历朝残部向中缅边地退缩} 的地点与行动范围属于京师与多地军政线路相接的尺度，行动者是清与南明永历朝残部。本节点回看「清军继续经营云南城镇与驿路，永历朝残部向中缅边地退缩」之前已可见的条件。占领昆明后清军仍须供应远征军、维持道路并争取地方中介，永历残部利用边地机动。读《清世祖实录》顺治十七年云南军报；《清史稿·吴三桂传》；内阁大库顺治十七年云南粮饷题本与《云南通志》时可见其记录边界，最小后果只能写作云南成为吴三桂等将领驻军和财政调度基地，也为后来的藩镇结构埋下条件。原有置信注记：清军控制主要城镇大体可靠；土司、矿区和山地实际控制不能以府城易手推定

\paragraph{康熙元年四月（1662年）：吴三桂军将自缅甸押回的永历帝朱由榔处死，南明朝廷名义中心结束}
\eventanchor{EQ-1662-YONGLI-EXECUTION}{吴三桂军将自缅甸押回的永历帝朱由榔处死，南明朝廷名义中心结束} 是清、南明永历朝与缅甸在边地道路、驻防与多方社群交错的尺度留下的编年标记。本节点叙述该事件本身。它从永历朝长期依赖流亡和边地补给，缅甸宫廷在清军压力下改变安置策略，吴三桂完成追捕走来；《清圣祖实录》康熙元年四月；《清史稿·永历帝本纪》；《琉璃宫史》缅文传统与《云南通志》固定可证部分，南明皇帝名义不复存在，但西南残部、郑氏与地方记忆仍继续影响清初整合构成下一步的可见结果。原有置信注记：被押回和处死主干可靠；缅甸移交过程、处死地点与责任叙事须比对中缅材料

\paragraph{康熙元年四月（1662年）：吴三桂军将自缅甸押回的永历帝朱由榔处死，南明朝廷名义中心结束}
本节点追踪「吴三桂军将自缅甸押回的永历帝朱由榔处死，南明朝廷名义中心结束」之后可见的余波：\eventanchor{EQ-1662-YONGLI-EXECUTION-AFTERMATH}{吴三桂军将自缅甸押回的永历帝朱由榔处死，南明朝廷名义中心结束}。清、南明永历朝与缅甸在边地道路、驻防与多方社群交错的尺度的处境可由永历朝长期依赖流亡和边地补给，缅甸宫廷在清军压力下改变安置策略，吴三桂完成追捕说明。《清圣祖实录》康熙元年四月；《清史稿·永历帝本纪》；《琉璃宫史》缅文传统与《云南通志》是本段材料入口，因而只能据此追踪到南明皇帝名义不复存在，但西南残部、郑氏与地方记忆仍继续影响清初整合。原有置信注记：被押回和处死主干可靠；缅甸移交过程、处死地点与责任叙事须比对中缅材料

\paragraph{康熙元年四月（1662年）：吴三桂军将自缅甸押回的永历帝朱由榔处死，南明朝廷名义中心结束}
\eventanchor{EQ-1662-YONGLI-EXECUTION-DECISION}{吴三桂军将自缅甸押回的永历帝朱由榔处死，南明朝廷名义中心结束} 发生在边地道路、驻防与多方社群交错的尺度，牵动清、南明永历朝与缅甸。本节点定位围绕「吴三桂军将自缅甸押回的永历帝朱由榔处死，南明朝廷名义中心结束」形成的处置与调度记录。永历朝长期依赖流亡和边地补给，缅甸宫廷在清军压力下改变安置策略，吴三桂完成追捕是其结构性条件；《清圣祖实录》康熙元年四月；《清史稿·永历帝本纪》；《琉璃宫史》缅文传统与《云南通志》留下的证据把读者带到南明皇帝名义不复存在，但西南残部、郑氏与地方记忆仍继续影响清初整合。原有置信注记：被押回和处死主干可靠；缅甸移交过程、处死地点与责任叙事须比对中缅材料

\paragraph{康熙元年四月（1662年）：吴三桂军将自缅甸押回的永历帝朱由榔处死，南明朝廷名义中心结束}
\eventanchor{EQ-1662-YONGLI-EXECUTION-PRELUDE}{吴三桂军将自缅甸押回的永历帝朱由榔处死，南明朝廷名义中心结束} 对应本节点回看「吴三桂军将自缅甸押回的永历帝朱由榔处死，南明朝廷名义中心结束」之前已可见的条件。在边地道路、驻防与多方社群交错的尺度，清、南明永历朝与缅甸无法绕开永历朝长期依赖流亡和边地补给，缅甸宫廷在清军压力下改变安置策略，吴三桂完成追捕。依据《清圣祖实录》康熙元年四月；《清史稿·永历帝本纪》；《琉璃宫史》缅文传统与《云南通志》，本段把后续影响限定为南明皇帝名义不复存在，但西南残部、郑氏与地方记忆仍继续影响清初整合。原有置信注记：被押回和处死主干可靠；缅甸移交过程、处死地点与责任叙事须比对中缅材料

\paragraph{康熙三年（1664年）：清军攻破夔东据点，李来亨等抗清武装失败，长江上游的明遗民军事中心瓦解}
\eventanchor{EQ-1664-KUIZHOU}{清军攻破夔东据点，李来亨等抗清武装失败，长江上游的明遗民军事中心瓦解}。本节点叙述该事件本身。清与夔东抗清武装在京师与多地军政线路相接的尺度的处境首先受永历朝覆灭后夔东武装仍凭山地、峡江与地方网络抵抗，清廷需清理长江航运安全隐患制约。《清圣祖实录》康熙三年湖广四川军报；《清史稿·李来亨传》；《夔州府志》与南明遗民文集留下可核的线索；可辨的后果是清廷加强川东鄂西通道控制，残余反清力量更难与外部政权形成连线。原有置信注记：夔东结局可靠；各部关系、山地据点与居民伤亡材料零散

\paragraph{康熙三年（1664年）：清军攻破夔东据点，李来亨等抗清武装失败，长江上游的明遗民军事中心瓦解}
\eventanchor{EQ-1664-KUIZHOU-AFTERMATH}{清军攻破夔东据点，李来亨等抗清武装失败，长江上游的明遗民军事中心瓦解} 置于京师与多地军政线路相接的尺度中阅读，本节点追踪「清军攻破夔东据点，李来亨等抗清武装失败，长江上游的明遗民军事中心瓦解」之后可见的余波。永历朝覆灭后夔东武装仍凭山地、峡江与地方网络抵抗，清廷需清理长江航运安全隐患构成行动者清与夔东抗清武装必须面对的条件。取证从《清圣祖实录》康熙三年湖广四川军报；《清史稿·李来亨传》；《夔州府志》与南明遗民文集出发，因而只能把变化谨慎地连到清廷加强川东鄂西通道控制，残余反清力量更难与外部政权形成连线。原有置信注记：夔东结局可靠；各部关系、山地据点与居民伤亡材料零散

\paragraph{康熙三年（1664年）：清军攻破夔东据点，李来亨等抗清武装失败，长江上游的明遗民军事中心瓦解}
\eventanchor{EQ-1664-KUIZHOU-DECISION}{清军攻破夔东据点，李来亨等抗清武装失败，长江上游的明遗民军事中心瓦解} 标记本节点定位围绕「清军攻破夔东据点，李来亨等抗清武装失败，长江上游的明遗民军事中心瓦解」形成的处置与调度记录，涉及清与夔东抗清武装。京师与多地军政线路相接的尺度不是抽象背景：永历朝覆灭后夔东武装仍凭山地、峡江与地方网络抵抗，清廷需清理长江航运安全隐患。《清圣祖实录》康熙三年湖广四川军报；《清史稿·李来亨传》；《夔州府志》与南明遗民文集所见的顺序随后指向清廷加强川东鄂西通道控制，残余反清力量更难与外部政权形成连线，但原有置信注记：夔东结局可靠；各部关系、山地据点与居民伤亡材料零散

\paragraph{康熙三年（1664年）：清军攻破夔东据点，李来亨等抗清武装失败，长江上游的明遗民军事中心瓦解}
在京师与多地军政线路相接的尺度，\eventanchor{EQ-1664-KUIZHOU-PRELUDE}{清军攻破夔东据点，李来亨等抗清武装失败，长江上游的明遗民军事中心瓦解} 让清与夔东抗清武装的一个选择或压力有了可查的坐标。本节点回看「清军攻破夔东据点，李来亨等抗清武装失败，长江上游的明遗民军事中心瓦解」之前已可见的条件。前因可写为永历朝覆灭后夔东武装仍凭山地、峡江与地方网络抵抗，清廷需清理长江航运安全隐患；《清圣祖实录》康熙三年湖广四川军报；《清史稿·李来亨传》；《夔州府志》与南明遗民文集限定了记录，后续变化是清廷加强川东鄂西通道控制，残余反清力量更难与外部政权形成连线。原有置信注记：夔东结局可靠；各部关系、山地据点与居民伤亡材料零散
